% !TEX program = xelatex
% Generated by docs/latex/build.py from the sources pinned in sources.json.
\documentclass[11pt,a4paper]{article}
\usepackage[a4paper,margin=25mm]{geometry}
\usepackage{fontspec}
\setmainfont{lmroman10-regular.otf}[BoldFont=lmroman10-bold.otf,ItalicFont=lmroman10-italic.otf,BoldItalicFont=lmroman10-bolditalic.otf]
\setsansfont{lmsans10-regular.otf}[BoldFont=lmsans10-bold.otf,ItalicFont=lmsans10-oblique.otf]
\setmonofont{JuliaMono-Regular.ttf}[Path=fonts/,Scale=0.83]
\usepackage{amsmath,amssymb,mathtools}
\usepackage{graphicx,calc,longtable,booktabs,array}
\usepackage{fvextra}
\usepackage{xurl}
\usepackage[colorlinks=true,linkcolor=blue,urlcolor=blue,citecolor=blue]{hyperref}
\providecommand{\tightlist}{\setlength{\itemsep}{0pt}\setlength{\parskip}{0pt}}
\setlength{\parindent}{0pt}
\setlength{\parskip}{6pt}
\setlength{\emergencystretch}{6em}
\begin{document}
\section{P1 --- Lean formalization appendix}\label{p1-lean-formalization-appendix}

This appendix prints the complete published Lean sources for P1 in the exact dependency order used by the package verifier. The source snapshot is Git commit \href{https://github.com/mpelteshki/climbing-to-the-frontier/commit/29fbb70da18564048f62492e43870f2f96a65463}{\texttt{29fbb70d\allowbreak{}a1856404\allowbreak{}8f62492e\allowbreak{}43870f2f\allowbreak{}96a65463}}, which matched the local \texttt{origin/main} tracking ref when this appendix was prepared. Every file below was read from that commit and checked byte for byte against the working copy and the recorded \href{https://github.com/mpelteshki/climbing-to-the-frontier/blob/29fbb70da18564048f62492e43870f2f96a65463/cagent/p1/evidence/source-sha256.txt}{source hashes}. Links to individual files are pinned to the same commit.

\subsection{Scope and proof boundary}\label{scope-and-proof-boundary}

The \href{https://github.com/mpelteshki/climbing-to-the-frontier/blob/29fbb70da18564048f62492e43870f2f96a65463/cagent/p1/README.md}{P1 package overview} and \href{https://github.com/mpelteshki/climbing-to-the-frontier/blob/29fbb70da18564048f62492e43870f2f96a65463/cagent/p1/writeup.md}{final written proof} mark C1--C5 \textbf{Solved by complete written arguments} under the task's written-proof hand-in rule. This appendix reports supporting Lean formalizations, \textbf{not} full Lean proofs of those five general statements. The package does not contain assembled general C1--C5 Lean theorems. C6 was skipped. No hackathon submission or organizer acceptance is claimed.

{\def\LTcaptype{none} % do not increment counter
\begin{longtable}[]{@{}
  >{\raggedright\arraybackslash}p{(\linewidth - 4\tabcolsep) * \real{0.3333}}
  >{\raggedright\arraybackslash}p{(\linewidth - 4\tabcolsep) * \real{0.3333}}
  >{\raggedright\arraybackslash}p{(\linewidth - 4\tabcolsep) * \real{0.3333}}@{}}
\toprule\noalign{}
\begin{minipage}[b]{\linewidth}\raggedright
Cell
\end{minipage} & \begin{minipage}[b]{\linewidth}\raggedright
Complete written result
\end{minipage} & \begin{minipage}[b]{\linewidth}\raggedright
Scope certified by these Lean modules
\end{minipage} \\
\midrule\noalign{}
\endhead
\bottomrule\noalign{}
\endlastfoot
C1 & All \(N\) planar lines, including empty and repeated configurations & \(N=2\) via \texttt{angle2\_le}, and named \(N=3,4\) theorems in \texttt{Angles.lean} \\
C2 & All \(m\ge2\) with non-neighbors orthogonal & \(m=2\) in \(\mathbb{R}^1\) and \(m=3\) in \(\mathbb{R}^2\) in \texttt{Angles.lean} \\
C3 & All \(d\ge1\) & \(d=1,2\) in \texttt{Angles.lean} \\
C4 & Five lines in \(\mathbb{R}^3\) and six in \(\mathbb{R}^4\), both sharp & Analytic and coordinate lemmas, compact maximizer, sparse degree-two reduction, corank bounds, and cycle inequalities; the final whole-cell inequalities are written proofs \\
C5 & All \(d\ge2\) with \(N=d+2\), sharp & \(d=2\) and shared C4 lemmas; the general induction is a written proof \\
C6 & Skipped & None \\
\end{longtable}
}

The \href{https://github.com/mpelteshki/climbing-to-the-frontier/blob/29fbb70da18564048f62492e43870f2f96a65463/cagent/p1/criteria-audit.md}{criteria audit} maps the complete written arguments to the live cell statements. The source files use neither \texttt{sorry}, \texttt{admit}, custom \texttt{axiom}, \texttt{unsafe}, nor \texttt{native\_decide} to assert completion. The \href{https://github.com/mpelteshki/climbing-to-the-frontier/blob/29fbb70da18564048f62492e43870f2f96a65463/cagent/p1/evidence/replay/README.md}{fresh local replay record} reports successful sequential compilation of these exact 18 files and 51 \texttt{\#print\ axioms} declarations; their reported axioms were only \texttt{propext}, \texttt{Classical.choice}, and \texttt{Quot.sound}. That replay reused the pinned Mathlib checkout; it was not a clean network download.

\subsection{Reproduction and exact source hashes}\label{reproduction-and-exact-source-hashes}

The package pins Lean \textbf{4.34.1} in \href{https://github.com/mpelteshki/climbing-to-the-frontier/blob/29fbb70da18564048f62492e43870f2f96a65463/cagent/p1/lean-toolchain}{\texttt{lean-toolchain}} and Mathlib commit \texttt{d13f23b7\allowbreak{}23b8a846\allowbreak{}827a245b\allowbreak{}89c10fc7\allowbreak{}d3f11612} in \href{https://github.com/mpelteshki/climbing-to-the-frontier/blob/29fbb70da18564048f62492e43870f2f96a65463/cagent/p1/lakefile.toml}{\texttt{lakefile.toml}}. From \texttt{cagent/p1} in the published repository, the portable replay command is:

\begin{Verbatim}[breaklines,breakanywhere,fontsize=\footnotesize]
./verify.sh
\end{Verbatim}

The \href{https://github.com/mpelteshki/climbing-to-the-frontier/blob/29fbb70da18564048f62492e43870f2f96a65463/cagent/p1/verify.sh}{verification script} fetches required Mathlib dependencies and compiles the following modules in order. The hashes are SHA-256 of each complete \texttt{.lean} file as read from the pinned Git commit:

\begin{Verbatim}[breaklines,breakanywhere,fontsize=\footnotesize]
Angles.lean a4f6c5b156b5b6630cb3c85bc0b342e3b32eb9957368a6cd7e79871c4bb019fe
C4Pentagon.lean 092e302e13dcc62784280594511a29a058f6b8191c56fac98ae09a18e33f964d
C4Extremum.lean 418cf7266191c243478469276cd84d95f054942949e478b8bc78d9eed9a33ba9
C4Rotation.lean af61ab1e5d9de2d496184a668e5bf66e8687b60d869f17ecf4ddf64685f8ae3a
C4SparseStep.lean 072592200ed4d8059b3a1538911a0e2c7dba50ec4ff738ea105db497cee02272
C4SignedSparse.lean 5bdcda7a554ed29ed6221743fe95a077ba52aba441239e0f6024bdc526db84a2
C4Geometry.lean d08fae0bc5dd30adc77f62237269fd36f9e0c1fb2c3c56d0a0d48c5cd014dc4f
C4LocalRotation.lean e87a6df1b1713b3192a22afcedbc00bf210a0be7f51b50a4c3e2f655060635a0
C4Selection.lean 768979c102b411b643bdf02eb47f8ef0b7535aeab151f745baaa00ae556e6720
C4Configuration.lean f43198c702add47aeb0477c1b85979a8b576e68b6c4cfd3035d6e155de3991d5
C4Replacement.lean 4d9286153ac302d4bf6d93e87418bffa41790380e034b574cc7667cd702222d3
C4Direction.lean a1b1b936ccf804d516739f861176d80e25368529f39af3cdd8d18c564005c953
C4MaximalSparse.lean 370d792a5e0d7b6b1b56dfbdf57577b68eac91b9690b8b328caaeb6019513218
C4NeighborCount.lean 9d4f9fa2f09f805192d2ac8b9d9a3e9801615b11b0f513fa83e0c1abb5d74833
C4SparseConfiguration.lean 7bd4a7aa9e3b8b48313688d6e6126ab219cf60c54ad275089fe8e9fd050f9932
C4MatrixCorank.lean 180d00723bd73798394275b59b01dfa85a2c797b335e8267dd121993b5224cee
C4FiveCycle.lean fe0bb40d9fc9b71f42d69ef814c7199db5a214bfbcee507addca3340f8d4aed7
C4SixCycle.lean 2dc09cb3ac16e054a07e888b0e9aa527f541f4537c84b906b34ec0646df1be6f
\end{Verbatim}

\subsection{Complete Lean source, in dependency order}\label{complete-lean-source-in-dependency-order}

\subsubsection{1. Angles.lean}\label{angles.lean}

\href{https://github.com/mpelteshki/climbing-to-the-frontier/blob/29fbb70da18564048f62492e43870f2f96a65463/cagent/p1/Angles.lean}{Pinned source file} · SHA-256 \texttt{a4f6c5b1\allowbreak{}56b5b663\allowbreak{}0cb3c85b\allowbreak{}c0b342e3\allowbreak{}b32eb995\allowbreak{}7368a6cd\allowbreak{}7e79871c\allowbreak{}4bb019fe}

\begin{Verbatim}[breaklines,breakanywhere,fontsize=\footnotesize]
import Mathlib.Analysis.SpecialFunctions.Trigonometric.Inverse
import Mathlib.Tactic.Linarith
import Mathlib.Tactic.Ring
import Mathlib.Tactic.NormNum

/-! Genuine Euclidean cases of P1. Vectors use Cartesian coordinates;
`unit2` is exactly squared Euclidean norm = 1. No general cell claimed here. -/
namespace Angles

abbrev Vec2 := ℝ × ℝ

def dot2 (x y : Vec2) : ℝ := x.1 * y.1 + x.2 * y.2
def unit2 (x : Vec2) : Prop := dot2 x x = 1
noncomputable def angle2 (x y : Vec2) : ℝ := Real.arccos |dot2 x y|

 theorem angle2_nonneg (x y : Vec2) : 0 ≤ angle2 x y := Real.arccos_nonneg _
 theorem angle2_le (x y : Vec2) : angle2 x y ≤ Real.pi / 2 :=
  Real.arccos_le_pi_div_two.mpr (abs_nonneg _)

 theorem dot2_bound {x y : Vec2} (hx : unit2 x) (hy : unit2 y) :
    |dot2 x y| ≤ 1 := by
  dsimp [unit2, dot2] at *
  have h := sq_nonneg (x.1 * y.2 - x.2 * y.1)
  have hid : (x.1*y.1+x.2*y.2)^2 + (x.1*y.2-x.2*y.1)^2 =
      (x.1*x.1+x.2*x.2)*(y.1*y.1+y.2*y.2) := by ring
  rw [hx, hy] at hid
  rw [abs_le]
  constructor <;> nlinarith [sq_nonneg (x.1*y.1+x.2*y.2+1),
    sq_nonneg (x.1*y.1+x.2*y.2-1)]

 theorem gram2 (x y z : Vec2) :
    dot2 x x * dot2 y y * dot2 z z + 2 * dot2 x y * dot2 y z * dot2 x z -
    dot2 x x * (dot2 y z)^2 - dot2 y y * (dot2 x z)^2 -
    dot2 z z * (dot2 x y)^2 = 0 := by
  dsimp [dot2]
  ring

 theorem arccos_pair {a b : ℝ} (ha : 0 ≤ a) (hb : 0 ≤ b)
    (h : a^2 + b^2 = 1) : Real.arccos a + Real.arccos b = Real.pi / 2 := by
  have hs : Real.sqrt (1 - a^2) = b := by
    rw [show 1-a^2 = b^2 by nlinarith, Real.sqrt_sq hb]
  rw [Real.arccos_eq_arcsin ha, hs, Real.arccos_eq_pi_div_two_sub_arcsin]
  ring

/-- C2, full geometric case m=3 (not arbitrary m). -/
 theorem c2_m3 {x y z : Vec2} (hx : unit2 x) (hy : unit2 y) (hz : unit2 z)
    (hxz : dot2 x z = 0) : angle2 x y + angle2 y z = Real.pi / 2 := by
  have hg := gram2 x y z
  change dot2 x x = 1 at hx
  change dot2 y y = 1 at hy
  change dot2 z z = 1 at hz
  rw [hx, hy, hz, hxz] at hg
  apply arccos_pair (abs_nonneg _) (abs_nonneg _)
  simpa only [sq_abs] using (show (dot2 x y)^2 + (dot2 y z)^2 = 1 by nlinarith [hg])

/-- Analytic certificate used for three arbitrary planar lines. -/
 theorem arccos_triple {a b c : ℝ} (ha : 0 ≤ a) (hb : 0 ≤ b) (hc : 0 ≤ c)
    (ha1 : a ≤ 1) (hb1 : b ≤ 1) (_hc1 : c ≤ 1)
    (h : 1 ≤ a^2 + b^2 + c^2 + 2*a*b*c) :
    Real.arccos a + Real.arccos b + Real.arccos c ≤ Real.pi := by
  have hsa : (Real.sqrt (1-a^2))^2 = 1-a^2 :=
    Real.sq_sqrt (by nlinarith)
  have hsb : (Real.sqrt (1-b^2))^2 = 1-b^2 :=
    Real.sq_sqrt (by nlinarith)
  have hsab : Real.sqrt (1-a^2) * Real.sqrt (1-b^2) ≤ c+a*b := by
    have hprod : (Real.sqrt (1-a^2) * Real.sqrt (1-b^2))^2 =
        (1-a^2)*(1-b^2) := by rw [mul_pow, hsa, hsb]
    have hnn := mul_nonneg ha hb
    have hsn := mul_nonneg (Real.sqrt_nonneg (1-a^2)) (Real.sqrt_nonneg (1-b^2))
    nlinarith [sq_nonneg (c+a*b-Real.sqrt (1-a^2)*Real.sqrt (1-b^2))]
  have hcosa : Real.cos (Real.arccos a) = a := Real.cos_arccos (by linarith) ha1
  have hcosb : Real.cos (Real.arccos b) = b := Real.cos_arccos (by linarith) hb1
  have hcos : -Real.cos (Real.arccos a + Real.arccos b) ≤ c := by
    rw [Real.cos_add, hcosa, hcosb, Real.sin_arccos, Real.sin_arccos]
    linarith
  have hsum0 : 0 ≤ Real.arccos a + Real.arccos b :=
    add_nonneg (Real.arccos_nonneg _) (Real.arccos_nonneg _)
  have hsump : Real.arccos a + Real.arccos b ≤ Real.pi := by
    linarith [Real.arccos_le_pi_div_two.mpr ha, Real.arccos_le_pi_div_two.mpr hb]
  have hc' := Real.arccos_le_arccos hcos
  rw [Real.arccos_neg, Real.arccos_cos hsum0 hsump] at hc'
  linarith

/-- C1 with N=3: all real unit vectors, no finite grid restriction. -/
 theorem c1_n3 {x y z : Vec2} (hx : unit2 x) (hy : unit2 y) (hz : unit2 z) :
    angle2 x y + angle2 x z + angle2 y z ≤ Real.pi := by
  have hg := gram2 x y z
  change dot2 x x = 1 at hx
  change dot2 y y = 1 at hy
  change dot2 z z = 1 at hz
  rw [hx, hy, hz] at hg
  have hp : -(dot2 x y * dot2 x z * dot2 y z) ≤
      |dot2 x y| * |dot2 x z| * |dot2 y z| := by
    calc
      _ ≤ |dot2 x y * dot2 x z * dot2 y z| := neg_le_abs _
      _ = _ := by simp only [abs_mul]
  apply arccos_triple (abs_nonneg _) (abs_nonneg _) (abs_nonneg _)
    (dot2_bound hx hy) (dot2_bound hx hz) (dot2_bound hy hz)
  simp only [sq_abs]
  nlinarith [hg, hp]

/-- C1, N=4, obtained by adding the four three-line inequalities. -/
 theorem c1_n4 {w x y z : Vec2} (hw : unit2 w) (hx : unit2 x)
    (hy : unit2 y) (hz : unit2 z) :
    angle2 w x + angle2 w y + angle2 w z + angle2 x y + angle2 x z + angle2 y z ≤
      2 * Real.pi := by
  linarith [c1_n3 hw hx hy, c1_n3 hw hx hz, c1_n3 hw hy hz, c1_n3 hx hy hz]

/-- C2, m=2, in R¹. This also supplies C3, d=1. -/
 theorem c2_m2 {x y : ℝ} (hx : x*x=1) (hy : y*y=1) :
    Real.arccos |x*y| = 0 := by
  have hs : (x*y)^2=1 := by
    calc
      (x*y)^2 = (x*x)*(y*y) := by ring
      _ = 1 := by rw [hx, hy]; norm_num
  have hab : |x*y|=1 := by
    have := sq_abs (x*y)
    have := abs_nonneg (x*y)
    nlinarith
  rw [hab, Real.arccos_one]

/-- C3, d=2, is the three-line planar case. -/
 theorem c3_d2 {x y z : Vec2} (hx : unit2 x) (hy : unit2 y) (hz : unit2 z) :
    angle2 x y + angle2 x z + angle2 y z ≤ Real.pi := c1_n3 hx hy hz

/-- C5, d=2, is the four-line planar case. -/
 theorem c5_d2 {w x y z : Vec2} (hw : unit2 w) (hx : unit2 x)
    (hy : unit2 y) (hz : unit2 z) :
    angle2 w x + angle2 w y + angle2 w z + angle2 x y + angle2 x z + angle2 y z ≤
      2 * Real.pi := c1_n4 hw hx hy hz

#print axioms c2_m2
#print axioms c1_n4
#print axioms c3_d2
#print axioms c5_d2
#print axioms c2_m3
#print axioms c1_n3
end Angles
\end{Verbatim}

\subsubsection{2. C4Pentagon.lean}\label{c4pentagon.lean}

\href{https://github.com/mpelteshki/climbing-to-the-frontier/blob/29fbb70da18564048f62492e43870f2f96a65463/cagent/p1/C4Pentagon.lean}{Pinned source file} · SHA-256 \texttt{092e302e\allowbreak{}13dcc627\allowbreak{}84280594\allowbreak{}511a29a0\allowbreak{}58f6b819\allowbreak{}1c56fac9\allowbreak{}8ae09a18\allowbreak{}e33f964d}

\begin{Verbatim}[breaklines,breakanywhere,fontsize=\footnotesize]
import Mathlib.Analysis.SpecialFunctions.Trigonometric.Inverse
import Mathlib.Tactic.Linarith
import Mathlib.Tactic.Ring
import Mathlib.Tactic.NormNum

/-! A scalar pentagon inequality for acute angles. -/
namespace C4Pentagon

private theorem sqrt_one_sub_sq_ge_one_sub {a : ℝ} (ha0 : 0 ≤ a) (ha1 : a ≤ 1) :
    1 - a ≤ Real.sqrt (1 - a ^ 2) := by
  have hsq : (Real.sqrt (1 - a ^ 2)) ^ 2 = 1 - a ^ 2 :=
    Real.sq_sqrt (by nlinarith)
  have hs := Real.sqrt_nonneg (1 - a ^ 2)
  nlinarith [mul_nonneg ha0 (sub_nonneg.mpr ha1)]

/-- The scalar estimate needed to close the acute C4 pentagon argument. -/
theorem pentagon_scalar {a b : ℝ} (ha0 : 0 ≤ a) (ha1 : a ≤ 1)
    (hb0 : 0 ≤ b) (hb1 : b ≤ 1) :
    Real.arccos (a * b) +
      Real.arcsin ((Real.sqrt (1 - a ^ 2) * Real.sqrt (1 - b ^ 2)) / (1 + a * b)) ≤
    Real.arccos a + Real.arccos b := by
  let s := Real.sqrt (1 - a ^ 2)
  let t := Real.sqrt (1 - b ^ 2)
  let r := s * t
  let u := a * b
  let d := 1 + u
  let L := Real.sqrt (1 - u ^ 2)
  have hs0 : 0 ≤ s := Real.sqrt_nonneg _
  have ht0 : 0 ≤ t := Real.sqrt_nonneg _
  have hs2 : s ^ 2 = 1 - a ^ 2 := Real.sq_sqrt (by nlinarith)
  have ht2 : t ^ 2 = 1 - b ^ 2 := Real.sq_sqrt (by nlinarith)
  have hr0 : 0 ≤ r := mul_nonneg hs0 ht0
  have hu0 : 0 ≤ u := mul_nonneg ha0 hb0
  have hu1 : u ≤ 1 := by nlinarith [mul_nonneg (sub_nonneg.mpr ha1) hb0]
  have hdpos : 0 < d := by dsimp [d]; linarith
  have hL0 : 0 ≤ L := Real.sqrt_nonneg _
  have hL2 : L ^ 2 = 1 - u ^ 2 := Real.sq_sqrt (by nlinarith)
  have hL1 : L ≤ 1 := by nlinarith
  have hr2 : r ^ 2 = (1 - a ^ 2) * (1 - b ^ 2) := by
    dsimp [r]
    rw [mul_pow, hs2, ht2]
  have hidentity : d ^ 2 - r ^ 2 = (a + b) ^ 2 := by
    dsimp [d, u]
    rw [hr2]
    ring
  have hab0 : 0 ≤ a + b := by linarith
  have hrd : r ≤ d := by nlinarith [sq_nonneg (a + b)]
  have hratio0 : 0 ≤ r / d := div_nonneg hr0 hdpos.le
  have hratio1 : r / d ≤ 1 := (div_le_one hdpos).mpr hrd
  have hcosE : Real.cos (Real.arcsin (r / d)) = (a + b) / d := by
    rw [Real.cos_arcsin]
    have hrad : 0 ≤ 1 - (r / d) ^ 2 := by nlinarith [sq_nonneg (r / d - 1)]
    have hsq : (Real.sqrt (1 - (r / d) ^ 2)) ^ 2 = ((a + b) / d) ^ 2 := by
      rw [Real.sq_sqrt hrad]
      calc
        1 - (r / d) ^ 2 = (d ^ 2 - r ^ 2) / d ^ 2 := by field_simp
        _ = ((a + b) / d) ^ 2 := by rw [hidentity, div_pow]
    nlinarith only [hsq, Real.sqrt_nonneg (1 - (r / d) ^ 2),
      div_nonneg hab0 hdpos.le]
  have hrbound : (1 - a) * (1 - b) ≤ r := by
    have hsa : 1 - a ≤ s := sqrt_one_sub_sq_ge_one_sub ha0 ha1
    have htb : 1 - b ≤ t := sqrt_one_sub_sq_ge_one_sub hb0 hb1
    have h1a : 0 ≤ 1 - a := by linarith
    have h1b : 0 ≤ 1 - b := by linarith
    exact (mul_le_mul hsa htb h1b hs0).trans_eq rfl
  have hcosc : Real.cos (Real.arccos u) = u := Real.cos_arccos (by linarith) hu1
  have hcospa : Real.cos (Real.arccos a) = a := Real.cos_arccos (by linarith) ha1
  have hcosqb : Real.cos (Real.arccos b) = b := Real.cos_arccos (by linarith) hb1
  have hsinc : Real.sin (Real.arccos u) = L := Real.sin_arccos u
  have hsinpa : Real.sin (Real.arccos a) = s := Real.sin_arccos a
  have hsinqb : Real.sin (Real.arccos b) = t := Real.sin_arccos b
  have hsinE : Real.sin (Real.arcsin (r / d)) = r / d :=
    Real.sin_arcsin (by linarith) hratio1
  have hcosineq :
      Real.cos (Real.arccos a + Real.arccos b) ≤
      Real.cos (Real.arccos u + Real.arcsin (r / d)) := by
    rw [Real.cos_add, Real.cos_add, hcospa, hcosqb, hsinpa, hsinqb,
      hcosc, hsinc, hcosE, hsinE]
    have hfactor : u * (1 - a) * (1 - b) ≤ u * r := by
      calc
        u * (1 - a) * (1 - b) = u * ((1 - a) * (1 - b)) := by ring
        _ ≤ u * r := mul_le_mul_of_nonneg_left hrbound hu0
    have hL : L * r ≤ r := by simpa using mul_le_mul_of_nonneg_right hL1 hr0
    have htarget : (u - r) * d ≤ u * (a + b) - L * r := by
      dsimp [d]
      nlinarith only [hfactor, hL, hu0, hr0]
    change u - r ≤ u * ((a + b) / d) - L * (r / d)
    calc
      u - r ≤ (u * (a + b) - L * r) / d := (le_div_iff₀ hdpos).2 htarget
      _ = u * ((a + b) / d) - L * (r / d) := by ring
  have hleft0 : 0 ≤ Real.arccos u + Real.arcsin (r / d) :=
    add_nonneg (Real.arccos_nonneg _) (Real.arcsin_nonneg.mpr hratio0)
  have hleftpi : Real.arccos u + Real.arcsin (r / d) ≤ Real.pi := by
    linarith [Real.arccos_le_pi_div_two.mpr hu0, Real.arcsin_le_pi_div_two (r / d)]
  have hright0 : 0 ≤ Real.arccos a + Real.arccos b :=
    add_nonneg (Real.arccos_nonneg _) (Real.arccos_nonneg _)
  have hrightpi : Real.arccos a + Real.arccos b ≤ Real.pi := by
    linarith [Real.arccos_le_pi_div_two.mpr ha0, Real.arccos_le_pi_div_two.mpr hb0]
  have hangle : Real.arccos u + Real.arcsin (r / d) ≤
      Real.arccos a + Real.arccos b := by
    by_contra hn
    have hstrict := Real.strictAntiOn_cos ⟨hright0, hrightpi⟩
      ⟨hleft0, hleftpi⟩ (lt_of_not_ge hn)
    exact (not_lt_of_ge hcosineq) hstrict
  simpa only [s, t, r, u, d] using hangle

/-- Angle form of `pentagon_scalar` for angles in the first quadrant. -/
theorem pentagon_scalar_angles {p q : ℝ} (hp0 : 0 ≤ p) (hp1 : p ≤ Real.pi / 2)
    (hq0 : 0 ≤ q) (hq1 : q ≤ Real.pi / 2) :
    Real.arccos (Real.cos p * Real.cos q) +
      Real.arcsin ((Real.sin p * Real.sin q) / (1 + Real.cos p * Real.cos q)) ≤
    p + q := by
  have hcp : 0 ≤ Real.cos p := Real.cos_nonneg_of_mem_Icc ⟨by linarith, hp1⟩
  have hcq : 0 ≤ Real.cos q := Real.cos_nonneg_of_mem_Icc ⟨by linarith, hq1⟩
  have hsp : 0 ≤ Real.sin p :=
    Real.sin_nonneg_of_nonneg_of_le_pi hp0 (by linarith [Real.pi_pos])
  have hsq : 0 ≤ Real.sin q :=
    Real.sin_nonneg_of_nonneg_of_le_pi hq0 (by linarith [Real.pi_pos])
  have hsp' : Real.sqrt (1 - Real.cos p ^ 2) = Real.sin p := by
    rw [show 1 - Real.cos p ^ 2 = Real.sin p ^ 2 by
      nlinarith [Real.sin_sq_add_cos_sq p], Real.sqrt_sq hsp]
  have hsq' : Real.sqrt (1 - Real.cos q ^ 2) = Real.sin q := by
    rw [show 1 - Real.cos q ^ 2 = Real.sin q ^ 2 by
      nlinarith [Real.sin_sq_add_cos_sq q], Real.sqrt_sq hsq]
  have h := pentagon_scalar hcp (Real.cos_le_one p) hcq (Real.cos_le_one q)
  rw [hsp', hsq', Real.arccos_cos hp0 (by linarith [Real.pi_pos]),
    Real.arccos_cos hq0 (by linarith [Real.pi_pos])] at h
  exact h

/-- The reverse three-angle certificate used for the six-cycle estimate. -/
theorem arccos_triple_ge_pi {x y z : ℝ} (hx0 : 0 ≤ x) (hx1 : x ≤ 1)
    (hy0 : 0 ≤ y) (hy1 : y ≤ 1) (hz0 : 0 ≤ z)
    (h : x ^ 2 + y ^ 2 + z ^ 2 + 2 * x * y * z ≤ 1) :
    Real.pi ≤ Real.arccos x + Real.arccos y + Real.arccos z := by
  have hsx : (Real.sqrt (1 - x ^ 2)) ^ 2 = 1 - x ^ 2 :=
    Real.sq_sqrt (by nlinarith)
  have hsy : (Real.sqrt (1 - y ^ 2)) ^ 2 = 1 - y ^ 2 :=
    Real.sq_sqrt (by nlinarith)
  have hprod : (Real.sqrt (1 - x ^ 2) * Real.sqrt (1 - y ^ 2)) ^ 2 =
      (1 - x ^ 2) * (1 - y ^ 2) := by rw [mul_pow, hsx, hsy]
  have hxy0 : 0 ≤ x * y + z := by positivity
  have hroot0 : 0 ≤ Real.sqrt (1 - x ^ 2) * Real.sqrt (1 - y ^ 2) := by
    positivity
  have hroot : x * y + z ≤ Real.sqrt (1 - x ^ 2) * Real.sqrt (1 - y ^ 2) := by
    nlinarith only [h, hprod, hxy0, hroot0]
  have hcos : Real.cos (Real.arccos x + Real.arccos y) ≤ -z := by
    rw [Real.cos_add, Real.cos_arccos (by linarith) hx1,
      Real.cos_arccos (by linarith) hy1, Real.sin_arccos, Real.sin_arccos]
    linarith
  have hsum0 : 0 ≤ Real.arccos x + Real.arccos y :=
    add_nonneg (Real.arccos_nonneg _) (Real.arccos_nonneg _)
  have hsump : Real.arccos x + Real.arccos y ≤ Real.pi := by
    linarith [Real.arccos_le_pi_div_two.mpr hx0, Real.arccos_le_pi_div_two.mpr hy0]
  have hc' := Real.arccos_le_arccos hcos
  rw [Real.arccos_neg, Real.arccos_cos hsum0 hsump] at hc'
  linarith

/-- A spherical right-triangle relation identifies the gap angle. -/
theorem gap_angle_add_projection {B U T : ℝ}
    (hB0 : 0 ≤ B) (hB1 : B ≤ Real.pi / 2)
    (hU0 : 0 ≤ U) (hU1 : U ≤ Real.pi / 2)
    (hT0 : 0 ≤ T) (hT1 : T ≤ Real.pi / 2)
    (hrel : Real.sin B * Real.sin T = Real.cos B * Real.sin U * Real.cos T) :
    Real.arcsin ((Real.sin B * Real.sin U) / (1 + Real.cos B * Real.cos U)) + T =
      Real.arccos (Real.cos T * Real.cos U) := by
  let a := Real.cos B
  let b := Real.cos U
  let r := Real.sin B * Real.sin U
  let d := 1 + a * b
  have ha0 : 0 ≤ a := Real.cos_nonneg_of_mem_Icc ⟨by linarith, hB1⟩
  have hb0 : 0 ≤ b := Real.cos_nonneg_of_mem_Icc ⟨by linarith, hU1⟩
  have ha1 : a ≤ 1 := Real.cos_le_one B
  have hb1 : b ≤ 1 := Real.cos_le_one U
  have hsB0 : 0 ≤ Real.sin B :=
    Real.sin_nonneg_of_nonneg_of_le_pi hB0 (by linarith [Real.pi_pos])
  have hsU0 : 0 ≤ Real.sin U :=
    Real.sin_nonneg_of_nonneg_of_le_pi hU0 (by linarith [Real.pi_pos])
  have hr0 : 0 ≤ r := mul_nonneg hsB0 hsU0
  have hdpos : 0 < d := by dsimp [d]; nlinarith [mul_nonneg ha0 hb0]
  have hBsq : (Real.sin B) ^ 2 = 1 - a ^ 2 := by
    dsimp [a]
    nlinarith [Real.sin_sq_add_cos_sq B]
  have hUsq : (Real.sin U) ^ 2 = 1 - b ^ 2 := by
    dsimp [b]
    nlinarith [Real.sin_sq_add_cos_sq U]
  have hr2 : r ^ 2 = (1 - a ^ 2) * (1 - b ^ 2) := by
    dsimp [r]
    rw [mul_pow, hBsq, hUsq]
  have hidentity : d ^ 2 - r ^ 2 = (a + b) ^ 2 := by
    dsimp [d]
    rw [hr2]
    ring
  have hrd : r ≤ d := by nlinarith [sq_nonneg (a + b)]
  have hratio0 : 0 ≤ r / d := div_nonneg hr0 hdpos.le
  have hratio1 : r / d ≤ 1 := (div_le_one hdpos).mpr hrd
  have hcosE : Real.cos (Real.arcsin (r / d)) = (a + b) / d := by
    rw [Real.cos_arcsin]
    have hrad : 0 ≤ 1 - (r / d) ^ 2 := by nlinarith [sq_nonneg (r / d - 1)]
    have hsq : (Real.sqrt (1 - (r / d) ^ 2)) ^ 2 = ((a + b) / d) ^ 2 := by
      rw [Real.sq_sqrt hrad]
      calc
        1 - (r / d) ^ 2 = (d ^ 2 - r ^ 2) / d ^ 2 := by field_simp
        _ = ((a + b) / d) ^ 2 := by rw [hidentity, div_pow]
    nlinarith only [hsq, Real.sqrt_nonneg (1 - (r / d) ^ 2),
      div_nonneg (add_nonneg ha0 hb0) hdpos.le]
  have hsinE : Real.sin (Real.arcsin (r / d)) = r / d :=
    Real.sin_arcsin (by linarith) hratio1
  have hcos : Real.cos (Real.arcsin (r / d) + T) = Real.cos T * b := by
    rw [Real.cos_add, hcosE, hsinE]
    have hrelMul := congrArg (fun v : ℝ => v * Real.sin U) hrel
    have hUtrig := congrArg (fun v : ℝ => a * Real.cos T * v)
      (Real.sin_sq_add_cos_sq U)
    calc
      (a + b) / d * Real.cos T - r / d * Real.sin T =
          (Real.cos T * (a + b) - Real.sin T * r) / d := by ring
      _ = Real.cos T * b := by
        apply (div_eq_iff (ne_of_gt hdpos)).2
        dsimp [d, r, a, b] at *
        nlinarith only [hrelMul, hUtrig]
  have hangle0 : 0 ≤ Real.arcsin (r / d) + T :=
    add_nonneg (Real.arcsin_nonneg.mpr hratio0) hT0
  have hanglepi : Real.arcsin (r / d) + T ≤ Real.pi := by
    linarith [Real.arcsin_le_pi_div_two (r / d)]
  have htarget := Real.arccos_cos hangle0 hanglepi
  rw [hcos] at htarget
  simpa only [r, d, a, b] using htarget.symm

/-- Companion scalar inequality for the C4 six-cycle projection. -/
theorem six_cycle_scalar {T U V : ℝ} (hT0 : 0 ≤ T) (hT1 : T ≤ Real.pi / 2)
    (hU0 : 0 ≤ U) (hU1 : U ≤ Real.pi / 2)
    (hV0 : 0 ≤ V) (hV1 : V ≤ Real.pi / 2) :
    Real.arcsin (Real.sin U * Real.sin V) ≤
      Real.arccos (Real.cos T * Real.cos U) +
        Real.arccos (Real.sin T * Real.cos V) - Real.pi / 2 := by
  let x := Real.cos T * Real.cos U
  let y := Real.sin T * Real.cos V
  let z := Real.sin U * Real.sin V
  have hcT0 : 0 ≤ Real.cos T := Real.cos_nonneg_of_mem_Icc ⟨by linarith, hT1⟩
  have hcU0 : 0 ≤ Real.cos U := Real.cos_nonneg_of_mem_Icc ⟨by linarith, hU1⟩
  have hcV0 : 0 ≤ Real.cos V := Real.cos_nonneg_of_mem_Icc ⟨by linarith, hV1⟩
  have hsT0 : 0 ≤ Real.sin T := Real.sin_nonneg_of_nonneg_of_le_pi hT0 (by linarith [Real.pi_pos])
  have hsU0 : 0 ≤ Real.sin U := Real.sin_nonneg_of_nonneg_of_le_pi hU0 (by linarith [Real.pi_pos])
  have hsV0 : 0 ≤ Real.sin V := Real.sin_nonneg_of_nonneg_of_le_pi hV0 (by linarith [Real.pi_pos])
  have hx0 : 0 ≤ x := mul_nonneg hcT0 hcU0
  have hy0 : 0 ≤ y := mul_nonneg hsT0 hcV0
  have hz0 : 0 ≤ z := mul_nonneg hsU0 hsV0
  have hx1 : x ≤ 1 := by nlinarith [mul_nonneg (sub_nonneg.mpr (Real.cos_le_one T)) hcU0, Real.cos_le_one U]
  have hy1 : y ≤ 1 := by nlinarith [mul_nonneg (sub_nonneg.mpr (Real.sin_le_one T)) hcV0, Real.cos_le_one V]
  have hz1 : z ≤ 1 := by nlinarith [mul_nonneg (sub_nonneg.mpr (Real.sin_le_one U)) hsV0, Real.sin_le_one V]
  have hT : (Real.cos T) ^ 2 + (Real.sin T) ^ 2 = 1 := by
    nlinarith [Real.sin_sq_add_cos_sq T]
  have hU : (Real.cos U) ^ 2 + (Real.sin U) ^ 2 = 1 := by
    nlinarith [Real.sin_sq_add_cos_sq U]
  have hV : (Real.cos V) ^ 2 + (Real.sin V) ^ 2 = 1 := by
    nlinarith [Real.sin_sq_add_cos_sq V]
  have hcoeff : (Real.sin T) ^ 2 + (Real.cos T) ^ 2 * (Real.sin U) ^ 2 =
      (Real.sin U) ^ 2 + (Real.sin T) ^ 2 * (Real.cos U) ^ 2 := by
    have hcT : (Real.cos T) ^ 2 = 1 - (Real.sin T) ^ 2 := by linarith
    have hcU : (Real.cos U) ^ 2 = 1 - (Real.sin U) ^ 2 := by linarith
    rw [hcT, hcU]
    ring
  let w := Real.cos T * Real.sin U * Real.cos V -
    Real.sin T * Real.sin V * Real.cos U
  have hidentity : x ^ 2 + y ^ 2 + z ^ 2 + 2 * x * y * z + w ^ 2 = 1 := by
    calc
      x ^ 2 + y ^ 2 + z ^ 2 + 2 * x * y * z + w ^ 2 =
          (Real.cos T) ^ 2 * (Real.cos U) ^ 2 +
          (Real.cos V) ^ 2 * ((Real.sin T) ^ 2 + (Real.cos T) ^ 2 * (Real.sin U) ^ 2) +
          (Real.sin V) ^ 2 * ((Real.sin U) ^ 2 + (Real.sin T) ^ 2 * (Real.cos U) ^ 2) := by
        dsimp [x, y, z, w]
        ring
      _ = (Real.cos T) ^ 2 * (Real.cos U) ^ 2 +
          ((Real.cos V) ^ 2 + (Real.sin V) ^ 2) *
            ((Real.sin U) ^ 2 + (Real.sin T) ^ 2 * (Real.cos U) ^ 2) := by
        rw [hcoeff]
        ring
      _ = (Real.cos T) ^ 2 * (Real.cos U) ^ 2 +
          ((Real.sin U) ^ 2 + (Real.sin T) ^ 2 * (Real.cos U) ^ 2) := by rw [hV]; ring
      _ = ((Real.cos T) ^ 2 + (Real.sin T) ^ 2) * (Real.cos U) ^ 2 +
          (Real.sin U) ^ 2 := by ring
      _ = 1 := by rw [hT, one_mul, hU]
  have hcert : x ^ 2 + y ^ 2 + z ^ 2 + 2 * x * y * z ≤ 1 := by
    nlinarith only [hidentity, sq_nonneg w]
  have htri := arccos_triple_ge_pi hx0 hx1 hy0 hy1 hz0 hcert
  rw [show Real.arccos z = Real.pi / 2 - Real.arcsin z from rfl] at htri
  dsimp [x, y, z] at htri ⊢
  linarith

/-- The local C4 projection estimate under the exact right-triangle relations. -/
theorem local_projection_inequality {T B C U V : ℝ}
    (hT0 : 0 ≤ T) (hT1 : T ≤ Real.pi / 2)
    (hB0 : 0 ≤ B) (hB1 : B ≤ Real.pi / 2)
    (hC0 : 0 ≤ C) (hC1 : C ≤ Real.pi / 2)
    (hU0 : 0 ≤ U) (hU1 : U ≤ Real.pi / 2)
    (hV0 : 0 ≤ V) (hV1 : V ≤ Real.pi / 2)
    (h1 : Real.sin B * Real.sin T = Real.cos B * Real.sin U * Real.cos T)
    (h2 : Real.sin C * Real.cos T = Real.cos C * Real.sin V * Real.sin T) :
    Real.arcsin (Real.sin U * Real.sin V) ≤
      (B + U - Real.arccos (Real.cos B * Real.cos U)) +
      (C + V - Real.arccos (Real.cos C * Real.cos V)) := by
  have hgap1 := gap_angle_add_projection hB0 hB1 hU0 hU1 hT0 hT1 h1
  have hT'0 : 0 ≤ Real.pi / 2 - T := by linarith
  have hT'1 : Real.pi / 2 - T ≤ Real.pi / 2 := by linarith
  have h2' : Real.sin C * Real.sin (Real.pi / 2 - T) =
      Real.cos C * Real.sin V * Real.cos (Real.pi / 2 - T) := by
    simpa only [Real.sin_pi_div_two_sub, Real.cos_pi_div_two_sub] using h2
  have hgap2 := gap_angle_add_projection hC0 hC1 hV0 hV1 hT'0 hT'1 h2'
  rw [Real.cos_pi_div_two_sub] at hgap2
  have hbound1 := pentagon_scalar_angles hB0 hB1 hU0 hU1
  have hbound2 := pentagon_scalar_angles hC0 hC1 hV0 hV1
  have hsix := six_cycle_scalar hT0 hT1 hU0 hU1 hV0 hV1
  linarith

#print axioms pentagon_scalar
#print axioms pentagon_scalar_angles
#print axioms arccos_triple_ge_pi
#print axioms six_cycle_scalar
#print axioms gap_angle_add_projection
#print axioms local_projection_inequality
end C4Pentagon
\end{Verbatim}

\subsubsection{3. C4Extremum.lean}\label{c4extremum.lean}

\href{https://github.com/mpelteshki/climbing-to-the-frontier/blob/29fbb70da18564048f62492e43870f2f96a65463/cagent/p1/C4Extremum.lean}{Pinned source file} · SHA-256 \texttt{418cf726\allowbreak{}6191c243\allowbreak{}47846927\allowbreak{}6cd84d95\allowbreak{}f0549429\allowbreak{}49e478b8\allowbreak{}bc78d9ee\allowbreak{}d9a33ba9}

\begin{Verbatim}[breaklines,breakanywhere,fontsize=\footnotesize]
import Mathlib.Analysis.Convex.Function
import Mathlib.Analysis.SpecialFunctions.Trigonometric.Inverse
import Mathlib.Tactic.FieldSimp
import Mathlib.Tactic.Linarith
import Mathlib.Tactic.Ring

namespace C4Extremum

/-- An interior maximum of a convex function forces both endpoint values to equal it. -/
theorem convex_max_endpoints {f : ℝ → ℝ} {l r : ℝ}
    (hl : l < 0) (hr : 0 < r) (hconv : ConvexOn ℝ (Set.Icc l r) f)
    (hmax : ∀ t ∈ Set.Icc l r, f t ≤ f 0) :
    f l = f 0 ∧ f r = f 0 := by
  have hd : 0 < r - l := by linarith
  let a : ℝ := r / (r - l)
  let b : ℝ := -l / (r - l)
  have ha : 0 < a := div_pos hr hd
  have hb : 0 < b := div_pos (neg_pos.mpr hl) hd
  have hab : a + b = 1 := by
    dsimp [a, b]
    rw [← add_div]
    have hnum : r + -l = r - l := by ring
    rw [hnum]
    exact div_self (ne_of_gt hd)
  have hzero : a • l + b • r = (0 : ℝ) := by
    dsimp [a, b]
    field_simp [ne_of_gt hd]
    ring
  have hlin : l ∈ Set.Icc l r := ⟨le_rfl, by linarith⟩
  have hrin : r ∈ Set.Icc l r := ⟨by linarith, le_rfl⟩
  have hweighted := hconv.2 hlin hrin ha.le hb.le hab
  rw [hzero] at hweighted
  simp only [smul_eq_mul] at hweighted
  have hfl : f l ≤ f 0 := hmax l hlin
  have hfr : f r ≤ f 0 := hmax r hrin
  have hsum : a * f 0 + b * f 0 = f 0 := by
    calc
      a * f 0 + b * f 0 = (a + b) * f 0 := by ring
      _ = f 0 := by rw [hab]; ring
  constructor
  · apply le_antisymm hfl
    by_contra h
    have hlt : f l < f 0 := lt_of_not_ge h
    have hltmul : a * f l < a * f 0 := mul_lt_mul_of_pos_left hlt ha
    have hlemul : b * f r ≤ b * f 0 := mul_le_mul_of_nonneg_left hfr hb.le
    linarith
  · apply le_antisymm hfr
    by_contra h
    have hlt : f r < f 0 := lt_of_not_ge h
    have hltmul : b * f r < b * f 0 := mul_lt_mul_of_pos_left hlt hb
    have hlemul : a * f l ≤ a * f 0 := mul_le_mul_of_nonneg_left hfl ha.le
    linarith

/-- Abstract endpoint contradiction used after a rotation is shown to preserve constraints. -/
theorem rotation_endpoint_contradiction {α : Type*} (rotation : ℝ → α)
    (objective : α → ℝ) (complexity : α → ℕ) {x : α} {l r : ℝ}
    (hl : l < 0) (hr : 0 < r)
    (hconv : ConvexOn ℝ (Set.Icc l r) (objective ∘ rotation))
    (hzero : rotation 0 = x)
    (hmax : ∀ y, objective y ≤ objective x)
    (hselect : ∀ y, objective y = objective x → complexity y ≤ complexity x)
    (hgain : complexity x < complexity (rotation l) ∨
      complexity x < complexity (rotation r)) : False := by
  have hmax' : ∀ t ∈ Set.Icc l r,
      (objective ∘ rotation) t ≤ (objective ∘ rotation) 0 := by
    intro t _
    simpa only [Function.comp_apply, hzero] using hmax (rotation t)
  obtain ⟨hle, hre⟩ := convex_max_endpoints hl hr hconv hmax'
  rcases hgain with hgain | hgain
  · have heq : objective (rotation l) = objective x := by
      simpa only [Function.comp_apply, hzero] using hle
    exact (not_lt_of_ge (hselect (rotation l) heq)) hgain
  · have heq : objective (rotation r) = objective x := by
      simpa only [Function.comp_apply, hzero] using hre
    exact (not_lt_of_ge (hselect (rotation r) heq)) hgain

#print axioms convex_max_endpoints
#print axioms rotation_endpoint_contradiction

end C4Extremum
\end{Verbatim}

\subsubsection{4. C4Rotation.lean}\label{c4rotation.lean}

\href{https://github.com/mpelteshki/climbing-to-the-frontier/blob/29fbb70da18564048f62492e43870f2f96a65463/cagent/p1/C4Rotation.lean}{Pinned source file} · SHA-256 \texttt{af61ab1e\allowbreak{}5d9de2d4\allowbreak{}96184a66\allowbreak{}8e5bf66e\allowbreak{}8687b60d\allowbreak{}869f17ec\allowbreak{}f4ddf646\allowbreak{}85f8ae3a}

\begin{Verbatim}[breaklines,breakanywhere,fontsize=\footnotesize]
import Mathlib.Analysis.Convex.Deriv
import Mathlib.Analysis.Normed.Module.Convex
import Mathlib.Analysis.SpecialFunctions.Trigonometric.InverseDeriv
import Mathlib.Analysis.SpecialFunctions.Sqrt
import Mathlib.Tactic.Linarith
import Mathlib.Tactic.Ring
import Mathlib.Tactic.NormNum
import Mathlib.Tactic.FieldSimp

namespace C4Rotation

private theorem inner_sq_lt_one {r t : ℝ} (hr0 : 0 ≤ r) (hr1 : r < 1) :
    (r * Real.cos t) ^ 2 < 1 := by
  have hc := Real.abs_cos_le_one t
  have hcsq : (Real.cos t) ^ 2 ≤ 1 := by
    nlinarith [mul_nonneg (abs_nonneg (Real.cos t)) (sub_nonneg.mpr hc),
      sq_abs (Real.cos t)]
  have hr2 : r ^ 2 < 1 := by nlinarith
  nlinarith [mul_nonneg (sq_nonneg r) (sub_nonneg.mpr hcsq)]

private theorem denominator_pos {r t : ℝ} (hr0 : 0 ≤ r) (hr1 : r < 1) :
    0 < Real.sqrt (1 - (r * Real.cos t) ^ 2) := by
  exact Real.sqrt_pos.2 (by linarith [inner_sq_lt_one hr0 hr1 (t := t)])

theorem hasDerivAt_angle {r t : ℝ} (hr0 : 0 ≤ r) (hr1 : r < 1) :
    HasDerivAt (fun x : ℝ => Real.arccos (r * Real.cos x))
      (r * Real.sin t / Real.sqrt (1 - (r * Real.cos t) ^ 2)) t := by
  have hsq := inner_sq_lt_one hr0 hr1 (t := t)
  have hlt : r * Real.cos t < 1 := by nlinarith
  have hgt : -1 < r * Real.cos t := by nlinarith
  have houter := Real.hasDerivAt_arccos (by linarith : r * Real.cos t ≠ -1)
    (by linarith : r * Real.cos t ≠ 1)
  have hinner : HasDerivAt (fun x : ℝ => r * Real.cos x) (r * -Real.sin t) t :=
    (Real.hasDerivAt_cos t).const_mul r
  convert houter.comp t hinner using 1
  · rfl
  · simp only [div_eq_mul_inv]
    ring

theorem hasDerivAt_slope {r t : ℝ} (hr0 : 0 ≤ r) (hr1 : r < 1) :
    HasDerivAt
      (fun x : ℝ => r * Real.sin x / Real.sqrt (1 - (r * Real.cos x) ^ 2))
      (r * (1 - r ^ 2) * Real.cos t /
        ((1 - (r * Real.cos t) ^ 2) * Real.sqrt (1 - (r * Real.cos t) ^ 2))) t := by
  have hcos : HasDerivAt (fun x : ℝ => r * Real.cos x) (r * -Real.sin t) t :=
    (Real.hasDerivAt_cos t).const_mul r
  have hD : HasDerivAt (fun x : ℝ => 1 - (r * Real.cos x) ^ 2)
      (2 * r ^ 2 * Real.cos t * Real.sin t) t := by
    convert (hcos.pow 2).const_sub 1 using 1; ring
  have hDpos : 0 < 1 - (r * Real.cos t) ^ 2 := by
    linarith [inner_sq_lt_one hr0 hr1 (t := t)]
  have hsqrt : HasDerivAt
      (fun x : ℝ => Real.sqrt (1 - (r * Real.cos x) ^ 2))
      ((1 / (2 * Real.sqrt (1 - (r * Real.cos t) ^ 2))) *
        (2 * r ^ 2 * Real.cos t * Real.sin t)) t := by
    convert (Real.hasDerivAt_sqrt hDpos.ne').comp t hD using 1; rfl
  have hnum : HasDerivAt (fun x : ℝ => r * Real.sin x)
      (r * Real.cos t) t := (Real.hasDerivAt_sin t).const_mul r
  have hsqrt_ne : Real.sqrt (1 - (r * Real.cos t) ^ 2) ≠ 0 :=
    (denominator_pos hr0 hr1 (t := t)).ne'
  convert hnum.div hsqrt hsqrt_ne using 1
  have hsqrt2 : (Real.sqrt (1 - (r * Real.cos t) ^ 2)) ^ 2 =
      1 - (r * Real.cos t) ^ 2 := Real.sq_sqrt hDpos.le
  have htrig := Real.sin_sq_add_cos_sq t
  field_simp
  ring_nf at hsqrt2
  have hden : (1 - r ^ 2 * Real.cos t ^ 2) *
      Real.sqrt (1 - r ^ 2 * Real.cos t ^ 2) =
      Real.sqrt (1 - r ^ 2 * Real.cos t ^ 2) ^ 3 := by
    calc
      _ = (Real.sqrt (1 - r ^ 2 * Real.cos t ^ 2)) ^ 2 *
          Real.sqrt (1 - r ^ 2 * Real.cos t ^ 2) := by rw [hsqrt2]
      _ = _ := by ring
  rw [hden]
  congr 1
  have hnumid : Real.sqrt (1 - r ^ 2 * Real.cos t ^ 2) ^ 2 -
      r ^ 2 * Real.sin t ^ 2 = 1 - r ^ 2 := by
    rw [hsqrt2]
    nlinarith [htrig]
  rw [hnumid]
  ring

private theorem convexOn_angle_lt_one {r : ℝ} (hr0 : 0 ≤ r) (hr1 : r < 1) :
    ConvexOn ℝ (Set.Icc (-(Real.pi / 2)) (Real.pi / 2))
      (fun t => Real.arccos (r * Real.cos t)) := by
  let D : Set ℝ := Set.Icc (-(Real.pi / 2)) (Real.pi / 2)
  have hD : Convex ℝ D := convex_Icc _ _
  have hcont : ContinuousOn (fun t => Real.arccos (r * Real.cos t)) D :=
    (Real.continuous_arccos.comp (continuous_const.mul Real.continuous_cos)).continuousOn
  apply convexOn_of_hasDerivWithinAt2_nonneg hD hcont
    (fun t _ => (hasDerivAt_angle hr0 hr1).hasDerivWithinAt)
    (fun t _ => (hasDerivAt_slope hr0 hr1).hasDerivWithinAt)
  intro t ht
  have htD : t ∈ D := interior_subset ht
  have hcos : 0 ≤ Real.cos t := Real.cos_nonneg_of_mem_Icc htD
  have hr2 : 0 ≤ 1 - r ^ 2 := by nlinarith
  have hdp : 0 < 1 - (r * Real.cos t) ^ 2 := by
    linarith [inner_sq_lt_one hr0 hr1 (t := t)]
  have hsp : 0 < Real.sqrt (1 - (r * Real.cos t) ^ 2) :=
    denominator_pos hr0 hr1
  exact div_nonneg (mul_nonneg (mul_nonneg hr0 hr2) hcos)
    (mul_nonneg hdp.le hsp.le)

theorem convexOn_angle {r : ℝ} (hr0 : 0 ≤ r) (hr1 : r ≤ 1) :
    ConvexOn ℝ (Set.Icc (-(Real.pi / 2)) (Real.pi / 2))
      (fun t => Real.arccos (r * Real.cos t)) := by
  by_cases hr : r < 1
  · exact convexOn_angle_lt_one hr0 hr
  have hreq : r = 1 := by linarith
  subst r
  have hnorm : ConvexOn ℝ (Set.Icc (-(Real.pi / 2)) (Real.pi / 2))
      (fun t : ℝ => |t|) := by
    have hnorm0 : ConvexOn ℝ (Set.Icc (-(Real.pi / 2)) (Real.pi / 2))
        (norm : ℝ → ℝ) := convexOn_norm (convex_Icc _ _)
    convert hnorm0 using 1
  apply hnorm.congr
  intro t ht
  rcases ht with ⟨hl, hu⟩
  simp only [one_mul]
  rcases le_total 0 t with hnonneg | hnonpos
  · rw [abs_of_nonneg hnonneg, Real.arccos_cos hnonneg]
    linarith [Real.pi_pos]
  · have hneg : 0 ≤ -t := by linarith
    rw [abs_of_nonpos hnonpos, ← Real.cos_neg, Real.arccos_cos hneg]
    linarith [Real.pi_pos]

#print axioms C4Rotation.convexOn_angle

end C4Rotation
\end{Verbatim}

\subsubsection{5. C4SparseStep.lean}\label{c4sparsestep.lean}

\href{https://github.com/mpelteshki/climbing-to-the-frontier/blob/29fbb70da18564048f62492e43870f2f96a65463/cagent/p1/C4SparseStep.lean}{Pinned source file} · SHA-256 \texttt{07259220\allowbreak{}0ed4d805\allowbreak{}9b3a1538\allowbreak{}911a0e2c\allowbreak{}7dba50ec\allowbreak{}4ff738ea\allowbreak{}105db497\allowbreak{}cee02272}

\begin{Verbatim}[breaklines,breakanywhere,fontsize=\footnotesize]
import Mathlib.Analysis.SpecialFunctions.Trigonometric.Inverse
import Mathlib.Analysis.Convex.Function
import Mathlib.Tactic.Linarith
import Mathlib.Tactic.Ring
import Mathlib.Tactic.FieldSimp
import Mathlib.Data.Finset.Max
import C4Extremum
import C4Rotation

namespace C4SparseStep

/-- A positive initial inner product determines a cosine phase whose positive
semicircle contains the initial configuration strictly in its interior. -/
theorem cosine_phase {a b : ℝ} (ha : 0 < a) (hab : a ^ 2 + b ^ 2 ≤ 1) :
    ∃ r p : ℝ, 0 < r ∧ r ≤ 1 ∧ -(Real.pi / 2) < p ∧ p < Real.pi / 2 ∧
      ∀ t : ℝ, a * Real.cos t + b * Real.sin t = r * Real.cos (t - p) := by
  let r := Real.sqrt (a ^ 2 + b ^ 2)
  have hrad : 0 < a ^ 2 + b ^ 2 := by positivity
  have hr : 0 < r := Real.sqrt_pos.mpr hrad
  have hr2 : r ^ 2 = a ^ 2 + b ^ 2 := Real.sq_sqrt hrad.le
  have hr1 : r ≤ 1 := by nlinarith
  have hblo : -r < b := by nlinarith [sq_pos_of_pos ha]
  have hbhi : b < r := by nlinarith [sq_pos_of_pos ha]
  have hratio0 : -1 < b / r := by
    apply (lt_div_iff₀ hr).2
    linarith
  have hratio1 : b / r < 1 := (div_lt_one hr).2 hbhi
  let p := Real.arcsin (b / r)
  have hp0 : -(Real.pi / 2) < p := Real.neg_pi_div_two_lt_arcsin.mpr hratio0
  have hp1 : p < Real.pi / 2 := Real.arcsin_lt_pi_div_two.mpr hratio1
  have hsin : Real.sin p = b / r := Real.sin_arcsin hratio0.le hratio1.le
  have hcos : Real.cos p = a / r := by
    have hcos0 : 0 ≤ Real.cos p := Real.cos_arcsin_nonneg _
    have hsquare : (Real.cos p) ^ 2 = (a / r) ^ 2 := by
      have hid := Real.sin_sq_add_cos_sq p
      rw [hsin] at hid
      have hrat : (a / r) ^ 2 + (b / r) ^ 2 = 1 := by
        field_simp
        nlinarith only [hr2]
      nlinarith only [hid, hrat]
    nlinarith only [hsquare, hcos0, div_pos ha hr]
  refine ⟨r, p, hr, hr1, hp0, hp1, ?_⟩
  intro t
  rw [Real.cos_sub, hcos, hsin]
  field_simp

/-- A finite nonempty family of positive cosine arcs has a common closed
interval around zero, and one arc vanishes at its left endpoint. -/
theorem common_positive_segment {ι : Type*} [Fintype ι] [Nonempty ι]
    (p : ι → ℝ) (hp : ∀ i, -(Real.pi / 2) < p i ∧ p i < Real.pi / 2) :
    ∃ l u : ℝ, l < 0 ∧ 0 < u ∧
      (∀ i t, t ∈ Set.Icc l u → t - p i ∈ Set.Icc (-(Real.pi / 2)) (Real.pi / 2)) ∧
      ∃ i, Real.cos (l - p i) = 0 := by
  obtain ⟨imax, _, hmax⟩ := Finset.exists_max_image Finset.univ p Finset.univ_nonempty
  obtain ⟨imin, _, hmin⟩ := Finset.exists_min_image Finset.univ p Finset.univ_nonempty
  refine ⟨p imax - Real.pi / 2, p imin + Real.pi / 2, by linarith [(hp imax).2],
    by linarith [(hp imin).1], ?_, ?_⟩
  · intro i t ht
    have hi0 := hmax i (Finset.mem_univ i)
    have hi1 := hmin i (Finset.mem_univ i)
    constructor <;> linarith [ht.1, ht.2]
  · refine ⟨imax, ?_⟩
    rw [show p imax - Real.pi / 2 - p imax = -(Real.pi / 2) by ring,
      Real.cos_neg, Real.cos_pi_div_two]

/-- The finite-family part of the first-zero rotation argument. The individual
arc convexity hypothesis is discharged by the analytic rotation theorem. -/
theorem phase_family_endpoint {ι : Type*} [Fintype ι] [Nonempty ι]
    (r p : ι → ℝ) (hr : ∀ i, 0 ≤ r i)
    (hp : ∀ i, -(Real.pi / 2) < p i ∧ p i < Real.pi / 2)
    (hconv : ∀ i, ConvexOn ℝ (Set.Icc (-(Real.pi / 2)) (Real.pi / 2))
      (fun t => Real.arccos (r i * Real.cos t)))
    (hmax : ∀ t : ℝ,
      (∑ i, Real.arccos |r i * Real.cos (t - p i)|) ≤
      ∑ i, Real.arccos |r i * Real.cos (0 - p i)|) :
    ∃ t : ℝ,
      (∑ i, Real.arccos |r i * Real.cos (t - p i)|) =
        (∑ i, Real.arccos |r i * Real.cos (0 - p i)|) ∧
      ∃ i, r i * Real.cos (t - p i) = 0 := by
  obtain ⟨l, u, hl, hu, hseg, izero, hzero⟩ := common_positive_segment p hp
  let f : ℝ → ℝ := fun t => ∑ i, Real.arccos |r i * Real.cos (t - p i)|
  have heach (i : ι) : ConvexOn ℝ (Set.Icc l u)
      (fun t => Real.arccos |r i * Real.cos (t - p i)|) := by
    have hsub : Set.Icc l u ⊆
        (fun t => -p i + t) ⁻¹' Set.Icc (-(Real.pi / 2)) (Real.pi / 2) := by
      intro t ht
      simpa only [Set.mem_preimage, sub_eq_add_neg, add_comm] using hseg i t ht
    have hshift := ((hconv i).translate_right (-p i)).subset hsub (convex_Icc l u)
    apply hshift.congr
    intro t ht
    have hnonneg : 0 ≤ r i * Real.cos (t - p i) :=
      mul_nonneg (hr i) (Real.cos_nonneg_of_mem_Icc (hseg i t ht))
    change Real.arccos (r i * Real.cos (-p i + t)) =
      Real.arccos |r i * Real.cos (t - p i)|
    rw [show -p i + t = t - p i by ring, abs_of_nonneg hnonneg]
  have hsum : ConvexOn ℝ (Set.Icc l u) f := by
    refine ⟨convex_Icc l u, ?_⟩
    intro x hx y hy a b ha hb hab
    change (∑ i, Real.arccos |r i * Real.cos (a * x + b * y - p i)|) ≤
      a * (∑ i, Real.arccos |r i * Real.cos (x - p i)|) +
        b * (∑ i, Real.arccos |r i * Real.cos (y - p i)|)
    calc
      _ ≤ ∑ i, (a * Real.arccos |r i * Real.cos (x - p i)| +
          b * Real.arccos |r i * Real.cos (y - p i)|) := by
        apply Finset.sum_le_sum
        intro i _
        exact (heach i).2 hx hy ha hb hab
      _ = _ := by rw [Finset.sum_add_distrib, Finset.mul_sum, Finset.mul_sum]
  have hend := (C4Extremum.convex_max_endpoints hl hu hsum (fun t _ => hmax t)).1
  exact ⟨l, hend, izero, by rw [hzero, mul_zero]⟩

/-- A maximal finite sum of line angles can be rotated until one positive
initial inner product vanishes, without changing the sum. -/
theorem positive_coefficient_endpoint {ι : Type*} [Fintype ι] [Nonempty ι]
    (a b : ι → ℝ) (ha : ∀ i, 0 < a i)
    (hab : ∀ i, (a i) ^ 2 + (b i) ^ 2 ≤ 1)
    (hmax : ∀ t : ℝ,
      (∑ i, Real.arccos |a i * Real.cos t + b i * Real.sin t|) ≤
      ∑ i, Real.arccos |a i|) :
    ∃ t : ℝ,
      (∑ i, Real.arccos |a i * Real.cos t + b i * Real.sin t|) =
        (∑ i, Real.arccos |a i|) ∧
      ∃ i, a i * Real.cos t + b i * Real.sin t = 0 := by
  classical
  choose r p hr0 hr1 hp0 hp1 heq using fun i => cosine_phase (ha i) (hab i)
  have hsum (t : ℝ) :
      (∑ i, Real.arccos |r i * Real.cos (t - p i)|) =
      ∑ i, Real.arccos |a i * Real.cos t + b i * Real.sin t| := by
    apply Finset.sum_congr rfl
    intro i _
    rw [← heq i t]
  have hsum0 : (∑ i, Real.arccos |r i * Real.cos (0 - p i)|) =
      ∑ i, Real.arccos |a i| := by
    simpa only [Real.cos_zero, Real.sin_zero, mul_one, mul_zero, add_zero] using hsum 0
  obtain ⟨t, ht, i, hi⟩ := phase_family_endpoint r p (fun i => (hr0 i).le)
    (fun i => ⟨hp0 i, hp1 i⟩) (fun i => C4Rotation.convexOn_angle (hr0 i).le (hr1 i))
    (fun t => by rw [hsum t, hsum0]; exact hmax t)
  refine ⟨t, ?_, i, ?_⟩
  · rwa [hsum t, hsum0] at ht
  · rwa [heq i t]

#print axioms cosine_phase
#print axioms common_positive_segment
#print axioms phase_family_endpoint
#print axioms positive_coefficient_endpoint
end C4SparseStep
\end{Verbatim}

\subsubsection{6. C4SignedSparse.lean}\label{c4signedsparse.lean}

\href{https://github.com/mpelteshki/climbing-to-the-frontier/blob/29fbb70da18564048f62492e43870f2f96a65463/cagent/p1/C4SignedSparse.lean}{Pinned source file} · SHA-256 \texttt{5bdcda7a\allowbreak{}554ed29e\allowbreak{}d6221743\allowbreak{}fe95a077\allowbreak{}ba52aba4\allowbreak{}41239e0f\allowbreak{}6024bdc5\allowbreak{}26db84a2}

\begin{Verbatim}[breaklines,breakanywhere,fontsize=\footnotesize]
import C4SparseStep

namespace C4SignedSparse

/-- A maximizing rotation can create a new zero among the nonzero initial
coefficients; coefficients initially zero remain identically zero. -/
theorem coefficient_endpoint {ι : Type*} [Fintype ι]
    (a b : ι → ℝ) (hne : ∃ i, a i ≠ 0)
    (hzero : ∀ i, a i = 0 → b i = 0)
    (hab : ∀ i, a i ^ 2 + b i ^ 2 ≤ 1)
    (hmax : ∀ t : ℝ,
      (∑ i, Real.arccos |a i * Real.cos t + b i * Real.sin t|) ≤
        ∑ i, Real.arccos |a i|) :
    ∃ t : ℝ,
      (∑ i, Real.arccos |a i * Real.cos t + b i * Real.sin t|) =
        ∑ i, Real.arccos |a i| ∧
      (∀ i, a i = 0 → a i * Real.cos t + b i * Real.sin t = 0) ∧
      ∃ i, a i ≠ 0 ∧ a i * Real.cos t + b i * Real.sin t = 0 := by
  classical
  let J := {i : ι // a i ≠ 0}
  let K := {i : ι // ¬a i ≠ 0}
  have hJ : Nonempty J := by
    obtain ⟨i, hi⟩ := hne
    exact ⟨⟨i, hi⟩⟩
  let A : J → ℝ := fun j => |a j.1|
  let B : J → ℝ := fun j => if 0 < a j.1 then b j.1 else -b j.1
  have hA (j : J) : 0 < A j := abs_pos.mpr j.2
  have hAB (j : J) : A j ^ 2 + B j ^ 2 ≤ 1 := by
    by_cases hp : 0 < a j.1
    · simpa [A, B, hp, abs_of_pos hp] using hab j.1
    · have hn : a j.1 < 0 := lt_of_le_of_ne (le_of_not_gt hp) j.2
      simpa [A, B, hp, abs_of_neg hn] using hab j.1
  have hterm (j : J) (t : ℝ) :
      |A j * Real.cos t + B j * Real.sin t| =
        |a j.1 * Real.cos t + b j.1 * Real.sin t| := by
    by_cases hp : 0 < a j.1
    · simp [A, B, hp, abs_of_pos hp]
    · have hn : a j.1 < 0 := lt_of_le_of_ne (le_of_not_gt hp) j.2
      simp only [A, B, hp, ↓reduceIte, abs_of_neg hn]
      calc
        |-a j.1 * Real.cos t + -b j.1 * Real.sin t| =
            |-(a j.1 * Real.cos t + b j.1 * Real.sin t)| := by congr 1; ring
        _ = _ := abs_neg _
  let F : ℝ → ι → ℝ := fun t i => Real.arccos |a i * Real.cos t + b i * Real.sin t|
  have hsplit (t : ℝ) :
      (∑ i, F t i) = (∑ j : J, F t j.1) + ∑ k : K, F t k.1 := by
    exact (Fintype.sum_subtype_add_sum_subtype (fun i => a i ≠ 0) (F t)).symm
  have hK (t : ℝ) : (∑ k : K, F t k.1) = ∑ k : K, F 0 k.1 := by
    apply Finset.sum_congr rfl
    intro k _
    have ha : a k.1 = 0 := by simpa using k.2
    simp [F, ha, hzero k.1 ha]
  have hmaxJ (t : ℝ) :
      (∑ j : J, Real.arccos |A j * Real.cos t + B j * Real.sin t|) ≤
        ∑ j : J, Real.arccos |A j| := by
    have htot : (∑ i, F t i) ≤ ∑ i, F 0 i := by
      simpa [F] using hmax t
    rw [hsplit t, hsplit 0, hK t] at htot
    have hJbound : (∑ j : J, F t j.1) ≤ ∑ j : J, F 0 j.1 := by linarith
    simpa [F, A, hterm] using hJbound
  have : Nonempty J := hJ
  obtain ⟨t, hsumJ, j, hj⟩ :=
    C4SparseStep.positive_coefficient_endpoint A B hA hAB hmaxJ
  refine ⟨t, ?_, ?_, j.1, j.2, ?_⟩
  · have hsumJ' : (∑ j : J, F t j.1) = ∑ j : J, F 0 j.1 := by
      simpa [F, A, hterm] using hsumJ
    calc
      (∑ i, Real.arccos |a i * Real.cos t + b i * Real.sin t|) =
          (∑ i, F t i) := rfl
      _ = (∑ j : J, F t j.1) + ∑ k : K, F t k.1 := hsplit t
      _ = (∑ j : J, F 0 j.1) + ∑ k : K, F 0 k.1 := by rw [hsumJ', hK]
      _ = ∑ i, F 0 i := (hsplit 0).symm
      _ = ∑ i, Real.arccos |a i| := by simp [F]
  · intro i hi
    simp [hi, hzero i hi]
  · by_cases hp : 0 < a j.1
    · simpa [A, B, hp, abs_of_pos hp] using hj
    · have hn : a j.1 < 0 := lt_of_le_of_ne (le_of_not_gt hp) j.2
      simp only [A, B, hp, ↓reduceIte, abs_of_neg hn] at hj
      nlinarith [hj]

#print axioms C4SignedSparse.coefficient_endpoint

end C4SignedSparse
\end{Verbatim}

\subsubsection{7. C4Geometry.lean}\label{c4geometry.lean}

\href{https://github.com/mpelteshki/climbing-to-the-frontier/blob/29fbb70da18564048f62492e43870f2f96a65463/cagent/p1/C4Geometry.lean}{Pinned source file} · SHA-256 \texttt{d08fae0b\allowbreak{}c5dd30ad\allowbreak{}c77f6223\allowbreak{}7269fd36\allowbreak{}f9e0c1fb\allowbreak{}2c3c56d0\allowbreak{}a0d48c5c\allowbreak{}d014dc4f}

\begin{Verbatim}[breaklines,breakanywhere,fontsize=\footnotesize]
import Mathlib.Analysis.InnerProductSpace.Basic
import Mathlib.Analysis.SpecialFunctions.Trigonometric.Inverse
import Mathlib.Tactic.Linarith
import Mathlib.Tactic.Ring

open scoped InnerProductSpace

namespace C4Geometry

variable {E : Type*} [NormedAddCommGroup E] [InnerProductSpace ℝ E]

/-- Rotation in the orthonormal plane spanned by `x` and `y`. -/
noncomputable def rotation (x y : E) (t : ℝ) : E := Real.cos t • x + Real.sin t • y

theorem inner_rotation (x y z : E) (t : ℝ) :
    ⟪rotation x y t, z⟫_ℝ = ⟪x, z⟫_ℝ * Real.cos t + ⟪y, z⟫_ℝ * Real.sin t := by
  simp only [rotation, inner_add_left, real_inner_smul_left]
  ring

theorem rotation_preserves_orthogonality (x y z : E) (t : ℝ)
    (hxz : ⟪x, z⟫_ℝ = 0) (hyz : ⟪y, z⟫_ℝ = 0) :
    ⟪rotation x y t, z⟫_ℝ = 0 := by
  rw [inner_rotation, hxz, hyz]
  ring

theorem rotation_unit (x y : E) (t : ℝ)
    (hx : ‖x‖ = 1) (hy : ‖y‖ = 1) (hxy : ⟪x, y⟫_ℝ = 0) :
    ‖rotation x y t‖ = 1 := by
  have hxx : ⟪x, x⟫_ℝ = 1 := by rw [real_inner_self_eq_norm_sq, hx]; norm_num
  have hyy : ⟪y, y⟫_ℝ = 1 := by rw [real_inner_self_eq_norm_sq, hy]; norm_num
  have hyx : ⟪y, x⟫_ℝ = 0 := by rw [real_inner_comm, hxy]
  have hrot : ⟪rotation x y t, rotation x y t⟫_ℝ =
      (Real.cos t) ^ 2 + (Real.sin t) ^ 2 := by
    simp only [rotation, inner_add_left, inner_add_right,
      real_inner_smul_left, real_inner_smul_right, hxx, hyy, hxy, hyx]
    ring
  have hn : ‖rotation x y t‖ ^ 2 = 1 := by
    rw [← real_inner_self_eq_norm_sq, hrot]
    nlinarith [Real.sin_sq_add_cos_sq t]
  nlinarith [norm_nonneg (rotation x y t)]

theorem bessel_two (x y z : E)
    (hx : ‖x‖ = 1) (hy : ‖y‖ = 1) (hxy : ⟪x, y⟫_ℝ = 0)
    (hz : ‖z‖ = 1) :
    ⟪x, z⟫_ℝ ^ 2 + ⟪y, z⟫_ℝ ^ 2 ≤ 1 := by
  let a : ℝ := ⟪x, z⟫_ℝ
  let b : ℝ := ⟪y, z⟫_ℝ
  let p : E := a • x + b • y
  have hxx : ⟪x, x⟫_ℝ = 1 := by rw [real_inner_self_eq_norm_sq, hx]; norm_num
  have hyy : ⟪y, y⟫_ℝ = 1 := by rw [real_inner_self_eq_norm_sq, hy]; norm_num
  have hyx : ⟪y, x⟫_ℝ = 0 := by rw [real_inner_comm, hxy]
  have hzx : ⟪z, x⟫_ℝ = a := by rw [real_inner_comm]
  have hzy : ⟪z, y⟫_ℝ = b := by rw [real_inner_comm]
  have hpp : ⟪p, p⟫_ℝ = a ^ 2 + b ^ 2 := by
    simp only [p, inner_add_left, inner_add_right,
      real_inner_smul_left, real_inner_smul_right, hxx, hyy, hxy, hyx]
    ring
  have hzp : ⟪z, p⟫_ℝ = a ^ 2 + b ^ 2 := by
    simp only [p, inner_add_right, real_inner_smul_right, hzx, hzy]
    ring
  have hpn : ‖p‖ ^ 2 = a ^ 2 + b ^ 2 := by
    rw [← real_inner_self_eq_norm_sq]
    exact hpp
  have hres := norm_sub_sq_real z p
  rw [hz, hzp, hpn] at hres
  have hres0 : 0 ≤ ‖z - p‖ ^ 2 := sq_nonneg _
  dsimp [a, b] at *
  nlinarith

#print axioms inner_rotation
#print axioms rotation_preserves_orthogonality
#print axioms rotation_unit
#print axioms bessel_two

end C4Geometry
\end{Verbatim}

\subsubsection{8. C4LocalRotation.lean}\label{c4localrotation.lean}

\href{https://github.com/mpelteshki/climbing-to-the-frontier/blob/29fbb70da18564048f62492e43870f2f96a65463/cagent/p1/C4LocalRotation.lean}{Pinned source file} · SHA-256 \texttt{e87a6df1\allowbreak{}b1713b31\allowbreak{}92a22afc\allowbreak{}edbc00bf\allowbreak{}210a0be7\allowbreak{}f51b50a4\allowbreak{}c3e2f655\allowbreak{}060635a0}

\begin{Verbatim}[breaklines,breakanywhere,fontsize=\footnotesize]
import C4Geometry
import C4SignedSparse

open scoped InnerProductSpace

namespace C4LocalRotation

variable {E : Type*} [NormedAddCommGroup E] [InnerProductSpace ℝ E]

/-- If a line can move in an orthonormal plane while preserving every old
orthogonality, a maximal sum of its angles to a finite family can be preserved
while creating a new orthogonality. -/
theorem add_orthogonality {ι : Type*} [Fintype ι]
    (x y : E) (z : ι → E)
    (hx : ‖x‖ = 1) (hy : ‖y‖ = 1) (hxy : ⟪x, y⟫_ℝ = 0)
    (hz : ∀ i, ‖z i‖ = 1)
    (hactive : ∃ i, ⟪x, z i⟫_ℝ ≠ 0)
    (hpreserve : ∀ i, ⟪x, z i⟫_ℝ = 0 → ⟪y, z i⟫_ℝ = 0)
    (hmax : ∀ t : ℝ,
      (∑ i, Real.arccos |⟪C4Geometry.rotation x y t, z i⟫_ℝ|) ≤
      ∑ i, Real.arccos |⟪x, z i⟫_ℝ|) :
    ∃ t : ℝ, ‖C4Geometry.rotation x y t‖ = 1 ∧
      (∑ i, Real.arccos |⟪C4Geometry.rotation x y t, z i⟫_ℝ|) =
        (∑ i, Real.arccos |⟪x, z i⟫_ℝ|) ∧
      (∀ i, ⟪x, z i⟫_ℝ = 0 → ⟪C4Geometry.rotation x y t, z i⟫_ℝ = 0) ∧
      ∃ i, ⟪x, z i⟫_ℝ ≠ 0 ∧ ⟪C4Geometry.rotation x y t, z i⟫_ℝ = 0 := by
  have hcoeff : ∀ i, ⟪x, z i⟫_ℝ ^ 2 + ⟪y, z i⟫_ℝ ^ 2 ≤ 1 :=
    fun i => C4Geometry.bessel_two x y (z i) hx hy hxy (hz i)
  have hmax' : ∀ t : ℝ,
      (∑ i, Real.arccos |⟪x, z i⟫_ℝ * Real.cos t + ⟪y, z i⟫_ℝ * Real.sin t|) ≤
        ∑ i, Real.arccos |⟪x, z i⟫_ℝ| := by
    intro t
    simpa only [C4Geometry.inner_rotation] using hmax t
  obtain ⟨t, ht, hold, i, hi, hnew⟩ := C4SignedSparse.coefficient_endpoint
    (fun i => ⟪x, z i⟫_ℝ) (fun i => ⟪y, z i⟫_ℝ)
    hactive hpreserve hcoeff hmax'
  refine ⟨t, C4Geometry.rotation_unit x y t hx hy hxy, ?_, ?_, i, hi, ?_⟩
  · simpa only [C4Geometry.inner_rotation] using ht
  · intro j hj
    simpa only [C4Geometry.inner_rotation] using hold j hj
  · simpa only [C4Geometry.inner_rotation] using hnew

#print axioms add_orthogonality
end C4LocalRotation
\end{Verbatim}

\subsubsection{9. C4Selection.lean}\label{c4selection.lean}

\href{https://github.com/mpelteshki/climbing-to-the-frontier/blob/29fbb70da18564048f62492e43870f2f96a65463/cagent/p1/C4Selection.lean}{Pinned source file} · SHA-256 \texttt{768979c1\allowbreak{}02b411b6\allowbreak{}43bdf02e\allowbreak{}b47f8ef0\allowbreak{}b7535aea\allowbreak{}b151f745\allowbreak{}baaa00ae\allowbreak{}556e6720}

\begin{Verbatim}[breaklines,breakanywhere,fontsize=\footnotesize]
import Mathlib.Topology.Order.Compact
import Mathlib.Order.Preorder.Finite
import Mathlib.Analysis.SpecialFunctions.Trigonometric.Inverse

namespace C4Selection

/-- Choose a score maximizer with greatest natural-number complexity among all ties. -/
theorem compact_lex_max {α : Type*} [TopologicalSpace α]
    {K : Set α} {S : α → ℝ} {C : α → ℕ} {M : ℕ}
    (hK : IsCompact K) (hne : K.Nonempty) (hS : ContinuousOn S K)
    (hC : ∀ x ∈ K, C x ≤ M) :
    ∃ x ∈ K, (∀ y ∈ K, S y ≤ S x) ∧
      (∀ y ∈ K, S y = S x → C y ≤ C x) := by
  obtain ⟨x₀, hx₀, hmax⟩ := hK.exists_isMaxOn hne hS
  let T : Set α := {y | y ∈ K ∧ S y = S x₀}
  have hTne : T.Nonempty := ⟨x₀, hx₀, rfl⟩
  have hfinite : (C '' T).Finite := by
    apply (Set.finite_Icc 0 M).subset
    rintro n ⟨y, hy, rfl⟩
    exact ⟨Nat.zero_le _, hC y hy.1⟩
  obtain ⟨x, hcx⟩ := Set.Finite.exists_maximalFor' C T hfinite hTne
  have hx : x ∈ K ∧ S x = S x₀ := hcx.1
  refine ⟨x, hx.1, ?_, ?_⟩
  · intro y hy
    rw [hx.2]
    exact hmax hy
  · intro y hy hscore
    have hyT : y ∈ T := ⟨hy, hscore.trans hx.2⟩
    rcases le_total (C y) (C x) with h | h
    · exact h
    · exact hcx.2 hyT h

#print axioms compact_lex_max

end C4Selection
\end{Verbatim}

\subsubsection{10. C4Configuration.lean}\label{c4configuration.lean}

\href{https://github.com/mpelteshki/climbing-to-the-frontier/blob/29fbb70da18564048f62492e43870f2f96a65463/cagent/p1/C4Configuration.lean}{Pinned source file} · SHA-256 \texttt{f43198c7\allowbreak{}02add47a\allowbreak{}eb0477c1\allowbreak{}b85979a8\allowbreak{}b576e68b\allowbreak{}6c4cfd30\allowbreak{}35d6e155\allowbreak{}de3991d5}

\begin{Verbatim}[breaklines,breakanywhere,fontsize=\footnotesize]
import C4Selection
import Mathlib.Analysis.InnerProductSpace.Continuous
import Mathlib.Analysis.SpecialFunctions.Trigonometric.Inverse
import Mathlib.Topology.MetricSpace.ProperSpace

open scoped InnerProductSpace

namespace C4Configuration

variable {E : Type*} [NormedAddCommGroup E] [InnerProductSpace ℝ E]

abbrev UnitVector (E : Type*) [NormedAddCommGroup E] :=
  Metric.sphere (0 : E) 1

abbrev Configuration (E : Type*) [NormedAddCommGroup E] (n : ℕ) :=
  Fin n → UnitVector E

/-- Each unordered pair of indices appears once. -/
def pairs (n : ℕ) : Finset (Fin n × Fin n) :=
  Finset.univ.filter (fun p => p.1 < p.2)

noncomputable def totalAngle {n : ℕ} (v : Configuration E n) : ℝ :=
  ∑ p ∈ pairs n, Real.arccos |⟪(v p.1).1, (v p.2).1⟫_ℝ|

noncomputable def orthogonalityCount {n : ℕ} (v : Configuration E n) : ℕ :=
  ((pairs n).filter (fun p => ⟪(v p.1).1, (v p.2).1⟫_ℝ = 0)).card

theorem continuous_totalAngle (n : ℕ) :
    Continuous (totalAngle (E := E) (n := n)) := by
  classical
  unfold totalAngle
  apply continuous_finsetSum
  intro p _
  have hfirst : Continuous (fun v : Configuration E n => (v p.1).1) :=
    continuous_subtype_val.comp (continuous_apply p.1)
  have hsecond : Continuous (fun v : Configuration E n => (v p.2).1) :=
    continuous_subtype_val.comp (continuous_apply p.2)
  exact Real.continuous_arccos.comp ((hfirst.inner hsecond).abs)

theorem orthogonalityCount_le {n : ℕ} (v : Configuration E n) :
    orthogonalityCount v ≤ (pairs n).card := by
  classical
  exact Finset.card_filter_le _ _

/-- A geometric angle maximizer with the most orthogonal pairs among all ties. -/
theorem exists_lex_max {n : ℕ} [ProperSpace E] (u : UnitVector E) :
    ∃ v : Configuration E n,
      (∀ w : Configuration E n, totalAngle w ≤ totalAngle v) ∧
      (∀ w : Configuration E n, totalAngle w = totalAngle v →
        orthogonalityCount w ≤ orthogonalityCount v) := by
  have hne : (Set.univ : Set (Configuration E n)).Nonempty :=
    ⟨fun _ => u, Set.mem_univ _⟩
  obtain ⟨v, _, hmax, hselect⟩ :=
    C4Selection.compact_lex_max (K := Set.univ)
      (S := totalAngle (E := E) (n := n)) (C := orthogonalityCount)
      (M := (pairs n).card) isCompact_univ hne
      (continuous_totalAngle n).continuousOn
      (by intro w _; exact orthogonalityCount_le w)
  refine ⟨v, ?_, ?_⟩
  · intro w
    exact hmax w (Set.mem_univ w)
  · intro w hw
    exact hselect w (Set.mem_univ w) hw

#print axioms continuous_totalAngle
#print axioms orthogonalityCount_le
#print axioms exists_lex_max

end C4Configuration
\end{Verbatim}

\subsubsection{11. C4Replacement.lean}\label{c4replacement.lean}

\href{https://github.com/mpelteshki/climbing-to-the-frontier/blob/29fbb70da18564048f62492e43870f2f96a65463/cagent/p1/C4Replacement.lean}{Pinned source file} · SHA-256 \texttt{4d928615\allowbreak{}3ac302d4\allowbreak{}bf6d93e8\allowbreak{}7418bffa\allowbreak{}41790380\allowbreak{}e034b574\allowbreak{}cc7667cd\allowbreak{}702222d3}

\begin{Verbatim}[breaklines,breakanywhere,fontsize=\footnotesize]
import C4Configuration

open scoped InnerProductSpace

namespace C4Replacement

open C4Configuration

variable {E : Type*} [NormedAddCommGroup E] [InnerProductSpace ℝ E]

def replacement {n : ℕ} (v : Configuration E n) (i : Fin n)
    (w : UnitVector E) : Configuration E n := Function.update v i w

def pairInner {n : ℕ} (v : Configuration E n) (p : Fin n × Fin n) : ℝ :=
  ⟪(v p.1).1, (v p.2).1⟫_ℝ

noncomputable def pairAngle {n : ℕ} (v : Configuration E n) (p : Fin n × Fin n) : ℝ :=
  Real.arccos |pairInner v p|

def incidentPairs (n : ℕ) (i : Fin n) : Finset (Fin n × Fin n) :=
  (pairs n).filter (fun p => p.1 = i ∨ p.2 = i)

noncomputable def incidentAngle {n : ℕ} (v : Configuration E n) (i : Fin n) : ℝ :=
  ∑ p ∈ incidentPairs n i, pairAngle v p

theorem pairInner_replacement_nonincident {n : ℕ} (v : Configuration E n)
    (i : Fin n) (w : UnitVector E) (p : Fin n × Fin n)
    (h₁ : p.1 ≠ i) (h₂ : p.2 ≠ i) :
    pairInner (replacement v i w) p = pairInner v p := by
  simp [pairInner, replacement, Function.update_of_ne h₁, Function.update_of_ne h₂]

theorem pairAngle_replacement_nonincident {n : ℕ} (v : Configuration E n)
    (i : Fin n) (w : UnitVector E) (p : Fin n × Fin n)
    (h₁ : p.1 ≠ i) (h₂ : p.2 ≠ i) :
    pairAngle (replacement v i w) p = pairAngle v p := by
  rw [pairAngle, pairAngle, pairInner_replacement_nonincident v i w p h₁ h₂]

/-- A replacement changes only the incident contribution to the total angle. -/
theorem totalAngle_replacement_balance {n : ℕ} (v : Configuration E n)
    (i : Fin n) (w : UnitVector E) :
    totalAngle (replacement v i w) + incidentAngle v i =
      totalAngle v + incidentAngle (replacement v i w) i := by
  classical
  let q : Fin n × Fin n → Prop := fun p => p.1 = i ∨ p.2 = i
  have hother :
      ∑ p ∈ (pairs n).filter (fun p => ¬ q p), pairAngle (replacement v i w) p =
      ∑ p ∈ (pairs n).filter (fun p => ¬ q p), pairAngle v p := by
    apply Finset.sum_congr rfl
    intro p hp
    have hp' : ¬ q p := (Finset.mem_filter.mp hp).2
    have h₁ : p.1 ≠ i := fun h => hp' (Or.inl h)
    have h₂ : p.2 ≠ i := fun h => hp' (Or.inr h)
    exact pairAngle_replacement_nonincident v i w p h₁ h₂
  have hsplit (u : Configuration E n) :
      totalAngle u = incidentAngle u i +
        ∑ p ∈ (pairs n).filter (fun p => ¬ q p), pairAngle u p := by
    simp only [totalAngle, incidentAngle, incidentPairs, pairAngle]
    exact (Finset.sum_filter_add_sum_filter_not (pairs n) q
      (fun p => Real.arccos |pairInner u p|)).symm
  rw [hsplit (replacement v i w), hsplit v, hother]
  ring

theorem totalAngle_replacement_eq_of_incident_eq {n : ℕ}
    (v : Configuration E n) (i : Fin n) (w : UnitVector E)
    (h : incidentAngle (replacement v i w) i = incidentAngle v i) :
    totalAngle (replacement v i w) = totalAngle v := by
  have hb := totalAngle_replacement_balance v i w
  rw [h] at hb
  exact add_right_cancel hb

theorem incidentAngle_le_of_totalAngle_max {n : ℕ}
    (v : Configuration E n) (i : Fin n) (w : UnitVector E)
    (hmax : ∀ u : Configuration E n, totalAngle u ≤ totalAngle v) :
    incidentAngle (replacement v i w) i ≤ incidentAngle v i := by
  have hb := totalAngle_replacement_balance v i w
  have hm := hmax (replacement v i w)
  linarith

/-- The uniquely ordered pair corresponding to two distinct indices. -/
def edge {n : ℕ} (i j : Fin n) : Fin n × Fin n :=
  if i < j then (i, j) else (j, i)

theorem edge_mem_pairs {n : ℕ} {i j : Fin n} (hij : j ≠ i) :
    edge i j ∈ pairs n := by
  have hne : i ≠ j := Ne.symm hij
  rcases lt_or_gt_of_ne hne with h | h
  · simp [edge, pairs, h]
  · have hnot : ¬ i < j := not_lt.mpr h.le
    simp [edge, pairs, hnot, h]

theorem pairInner_edge {n : ℕ} (v : Configuration E n)
    (i j : Fin n) :
    pairInner v (edge i j) = ⟪(v i).1, (v j).1⟫_ℝ := by
  by_cases h : i < j
  · simp [pairInner, edge, h]
  · simp [pairInner, edge, h, real_inner_comm]

theorem pairInner_replacement_edge {n : ℕ} (v : Configuration E n)
    (i j : Fin n) (w : UnitVector E) (hji : j ≠ i) :
    pairInner (replacement v i w) (edge i j) = ⟪w.1, (v j).1⟫_ℝ := by
  by_cases h : i < j
  · simp [pairInner, edge, h, replacement, hji]
  · simp [pairInner, edge, h, replacement, hji, real_inner_comm]

theorem edge_injective_away {n : ℕ} (i j k : Fin n)
    (hji : j ≠ i) (hki : k ≠ i)
    (h : edge i j = edge i k) : j = k := by
  by_cases hj : i < j
  · by_cases hk : i < k
    · simpa [edge, hj, hk] using congrArg Prod.snd h
    · have hki' : k < i := (lt_or_gt_of_ne (Ne.symm hki)).resolve_left hk
      have hi : i = k := by simpa [edge, hj, hk] using congrArg Prod.fst h
      exact (hki hi.symm).elim
  · by_cases hk : i < k
    · have hji' : j < i := (lt_or_gt_of_ne (Ne.symm hji)).resolve_left hj
      have hi : j = i := by simpa [edge, hj, hk] using congrArg Prod.fst h
      exact (hji hi).elim
    · simpa [edge, hj, hk] using congrArg Prod.fst h

/-- The incident pair sum is the usual star sum over all other vertices. -/
theorem incidentAngle_eq_star {n : ℕ} (v : Configuration E n) (i : Fin n) :
    incidentAngle v i =
      ∑ j ∈ Finset.univ.erase i, Real.arccos |⟪(v i).1, (v j).1⟫_ℝ| := by
  classical
  unfold incidentAngle
  symm
  apply Finset.sum_bij (fun j _ => edge i j)
  · intro j hj
    have hji : j ≠ i := (Finset.mem_erase.mp hj).1
    apply Finset.mem_filter.mpr
    refine ⟨edge_mem_pairs hji, ?_⟩
    by_cases h : i < j
    · simp [edge, h]
    · simp [edge, h]
  · intro j hj k hk heq
    exact edge_injective_away i j k (Finset.mem_erase.mp hj).1
      (Finset.mem_erase.mp hk).1 heq
  · intro p hp
    obtain ⟨hpair, hinc⟩ := Finset.mem_filter.mp hp
    have hlt : p.1 < p.2 := (Finset.mem_filter.mp hpair).2
    rcases hinc with hleft | hright
    · refine ⟨p.2, Finset.mem_erase.mpr ⟨?_, Finset.mem_univ _⟩, ?_⟩
      · intro heq
        exact (ne_of_lt hlt) (hleft.trans heq.symm)
      · rcases p with ⟨a, b⟩
        simp only at hleft hlt
        subst a
        simp [edge, hlt]
    · refine ⟨p.1, Finset.mem_erase.mpr ⟨?_, Finset.mem_univ _⟩, ?_⟩
      · intro heq
        exact (ne_of_lt hlt) (heq.trans hright.symm)
      · rcases p with ⟨a, b⟩
        simp only at hright hlt
        subst b
        have hnot : ¬ i < a := not_lt.mpr hlt.le
        simp [edge, hnot]
  · intro j hj
    rw [pairAngle, pairInner_edge]

/-- The incident sum indexed by neighbors as a subtype. -/
theorem incidentAngle_eq_sum_neighbors {n : ℕ}
    (v : Configuration E n) (i : Fin n) :
    incidentAngle v i =
      ∑ j : {j : Fin n // j ≠ i},
        Real.arccos |⟪(v i).1, (v j.1).1⟫_ℝ| := by
  classical
  rw [incidentAngle_eq_star]
  symm
  apply Finset.sum_bij (fun j _ => j.1)
  · intro j _
    exact Finset.mem_erase.mpr ⟨j.property, Finset.mem_univ _⟩
  · intro j _ k _ h
    exact Subtype.ext h
  · intro j hj
    exact ⟨⟨j, (Finset.mem_erase.mp hj).1⟩, Finset.mem_univ _, rfl⟩
  · intro j _
    rfl

/-- Exact replacement balance with the usual neighbor-index sums. -/
theorem totalAngle_replacement_star_balance {n : ℕ}
    (v : Configuration E n) (i : Fin n) (w : UnitVector E) :
    totalAngle (replacement v i w) +
      (∑ j ∈ Finset.univ.erase i,
        Real.arccos |⟪(v i).1, (v j).1⟫_ℝ|) =
    totalAngle v +
      (∑ j ∈ Finset.univ.erase i,
        Real.arccos |⟪w.1, (v j).1⟫_ℝ|) := by
  classical
  have hb := totalAngle_replacement_balance v i w
  rw [incidentAngle_eq_star v i,
    incidentAngle_eq_star (replacement v i w) i] at hb
  have hnewsum :
      (∑ j ∈ Finset.univ.erase i,
        Real.arccos |⟪(replacement v i w i).1,
          (replacement v i w j).1⟫_ℝ|) =
      (∑ j ∈ Finset.univ.erase i,
        Real.arccos |⟪w.1, (v j).1⟫_ℝ|) := by
    apply Finset.sum_congr rfl
    intro j hj
    have hji : j ≠ i := (Finset.mem_erase.mp hj).1
    simp [replacement, hji]
  rw [hnewsum] at hb
  exact hb

theorem neighbor_sum_le_of_totalAngle_max {n : ℕ}
    (v : Configuration E n) (i : Fin n) (w : UnitVector E)
    (hmax : ∀ u : Configuration E n, totalAngle u ≤ totalAngle v) :
    (∑ j : {j : Fin n // j ≠ i},
      Real.arccos |⟪w.1, (v j.1).1⟫_ℝ|) ≤
    (∑ j : {j : Fin n // j ≠ i},
      Real.arccos |⟪(v i).1, (v j.1).1⟫_ℝ|) := by
  have h := incidentAngle_le_of_totalAngle_max v i w hmax
  rw [incidentAngle_eq_sum_neighbors (replacement v i w) i,
    incidentAngle_eq_sum_neighbors v i] at h
  have hnewsum :
      (∑ j : {j : Fin n // j ≠ i},
        Real.arccos |⟪(replacement v i w i).1,
          (replacement v i w j.1).1⟫_ℝ|) =
      (∑ j : {j : Fin n // j ≠ i},
        Real.arccos |⟪w.1, (v j.1).1⟫_ℝ|) := by
    apply Finset.sum_congr rfl
    intro j _
    simp [replacement, j.property]
  rw [hnewsum] at h
  exact h

theorem pairInner_zero_preserved {n : ℕ} (v : Configuration E n)
    (i : Fin n) (w : UnitVector E)
    (hpres : ∀ j : Fin n, j ≠ i →
      ⟪(v i).1, (v j).1⟫_ℝ = 0 → ⟪w.1, (v j).1⟫_ℝ = 0)
    (p : Fin n × Fin n) (hp : p ∈ pairs n)
    (hzero : pairInner v p = 0) :
    pairInner (replacement v i w) p = 0 := by
  by_cases hleft : p.1 = i
  · have hright : p.2 ≠ i := by
      intro heq
      have hlt : p.1 < p.2 := (Finset.mem_filter.mp hp).2
      simp [hleft, heq] at hlt
    have hold : ⟪(v i).1, (v p.2).1⟫_ℝ = 0 := by
      simpa [pairInner, hleft] using hzero
    have hnew := hpres p.2 hright hold
    simpa [pairInner, replacement, hleft, hright] using hnew
  · by_cases hright : p.2 = i
    · have hold : ⟪(v i).1, (v p.1).1⟫_ℝ = 0 := by
        rw [real_inner_comm]
        simpa [pairInner, hright] using hzero
      have hnew := hpres p.1 hleft hold
      rw [real_inner_comm] at hnew
      simpa [pairInner, replacement, hleft, hright] using hnew
    · rw [pairInner_replacement_nonincident v i w p hleft hright]
      exact hzero

/-- Preserving old incident zeros and creating one new zero increases the count. -/
theorem orthogonalityCount_replacement_strict {n : ℕ}
    (v : Configuration E n) (i : Fin n) (w : UnitVector E)
    (hpres : ∀ j : Fin n, j ≠ i →
      ⟪(v i).1, (v j).1⟫_ℝ = 0 → ⟪w.1, (v j).1⟫_ℝ = 0)
    (hgain : ∃ j : Fin n, j ≠ i ∧
      ⟪(v i).1, (v j).1⟫_ℝ ≠ 0 ∧ ⟪w.1, (v j).1⟫_ℝ = 0) :
    orthogonalityCount v < orthogonalityCount (replacement v i w) := by
  classical
  let old := (pairs n).filter (fun p => pairInner v p = 0)
  let new := (pairs n).filter (fun p => pairInner (replacement v i w) p = 0)
  have hsubset : old ⊆ new := by
    intro p hp
    obtain ⟨hpair, hzero⟩ := Finset.mem_filter.mp hp
    exact Finset.mem_filter.mpr
      ⟨hpair, pairInner_zero_preserved v i w hpres p hpair hzero⟩
  obtain ⟨j, hji, hold, hnew⟩ := hgain
  have hp : edge i j ∈ pairs n := edge_mem_pairs hji
  have hpnew : edge i j ∈ new := by
    apply Finset.mem_filter.mpr
    constructor
    · exact hp
    · rw [pairInner_replacement_edge v i j w hji]
      exact hnew
  have hpnot : edge i j ∉ old := by
    intro h
    have hz : pairInner v (edge i j) = 0 := (Finset.mem_filter.mp h).2
    rw [pairInner_edge] at hz
    exact hold hz
  have hstrict : old ⊂ new :=
    (Finset.ssubset_iff_of_subset hsubset).2 ⟨edge i j, hpnew, hpnot⟩
  exact Finset.card_lt_card hstrict

#print axioms totalAngle_replacement_balance
#print axioms totalAngle_replacement_eq_of_incident_eq
#print axioms incidentAngle_le_of_totalAngle_max
#print axioms incidentAngle_eq_star
#print axioms incidentAngle_eq_sum_neighbors
#print axioms totalAngle_replacement_star_balance
#print axioms neighbor_sum_le_of_totalAngle_max
#print axioms orthogonalityCount_replacement_strict

end C4Replacement
\end{Verbatim}

\subsubsection{12. C4Direction.lean}\label{c4direction.lean}

\href{https://github.com/mpelteshki/climbing-to-the-frontier/blob/29fbb70da18564048f62492e43870f2f96a65463/cagent/p1/C4Direction.lean}{Pinned source file} · SHA-256 \texttt{a1b1b936\allowbreak{}ccf804d5\allowbreak{}16739f86\allowbreak{}1176d80e\allowbreak{}25368529\allowbreak{}f39af3cd\allowbreak{}d8d18c56\allowbreak{}4005c953}

\begin{Verbatim}[breaklines,breakanywhere,fontsize=\footnotesize]
import Mathlib.Analysis.InnerProductSpace.Projection.FiniteDimensional
import Mathlib.LinearAlgebra.FiniteDimensional.Lemmas

open scoped InnerProductSpace

namespace C4Direction

/-- If a subspace and one extra vector leave room in a finite-dimensional
real inner-product space, choose a unit vector orthogonal to both. -/
theorem exists_unit_orthogonal {E : Type*} [NormedAddCommGroup E]
    [InnerProductSpace ℝ E] [FiniteDimensional ℝ E]
    (W : Submodule ℝ E) (x : E)
    (hdim : Module.finrank ℝ W + 1 < Module.finrank ℝ E) :
    ∃ y : E, ‖y‖ = 1 ∧ ⟪x, y⟫_ℝ = 0 ∧
      ∀ w ∈ W, ⟪y, w⟫_ℝ = 0 := by
  let K : Submodule ℝ E := W ⊔ ℝ ∙ x
  have hspan : Module.finrank ℝ (ℝ ∙ x) ≤ 1 := by
    simpa using (finrank_span_le_card ({x} : Set E))
  have hK : Module.finrank ℝ K ≤ Module.finrank ℝ W + 1 := by
    have hsup := Submodule.finrank_sup_add_finrank_inf_eq W (ℝ ∙ x)
    dsimp [K]
    omega
  have horth : 0 < Module.finrank ℝ Kᗮ := by
    have hdimorth := K.finrank_add_finrank_orthogonal
    omega
  have hne : Kᗮ ≠ ⊥ := by
    intro heq
    rw [heq, finrank_bot] at horth
    omega
  obtain ⟨z, hz, hz0⟩ := Submodule.exists_mem_ne_zero_of_ne_bot hne
  let y : E := (‖z‖⁻¹ : ℝ) • z
  have hy : y ∈ Kᗮ := Submodule.smul_mem _ _ hz
  have hynorm : ‖y‖ = 1 := by
    change ‖(‖z‖⁻¹ : ℝ) • z‖ = 1
    convert norm_smul_inv_norm (𝕜 := ℝ) hz0 using 1
  refine ⟨y, hynorm, ?_, ?_⟩
  · have hxK : x ∈ K := Submodule.mem_sup_right (Submodule.mem_span_singleton_self x)
    have h := (K.mem_orthogonal' y).mp hy x hxK
    rwa [real_inner_comm] at h
  · intro w hw
    exact (K.mem_orthogonal' y).mp hy w (Submodule.mem_sup_left hw)

#print axioms C4Direction.exists_unit_orthogonal

end C4Direction
\end{Verbatim}

\subsubsection{13. C4MaximalSparse.lean}\label{c4maximalsparse.lean}

\href{https://github.com/mpelteshki/climbing-to-the-frontier/blob/29fbb70da18564048f62492e43870f2f96a65463/cagent/p1/C4MaximalSparse.lean}{Pinned source file} · SHA-256 \texttt{370d792a\allowbreak{}5e0d7b6b\allowbreak{}1b56dfbd\allowbreak{}f57577b6\allowbreak{}8eac91b9\allowbreak{}690b8b32\allowbreak{}8caaeb60\allowbreak{}19513218}

\begin{Verbatim}[breaklines,breakanywhere,fontsize=\footnotesize]
import C4Configuration
import C4LocalRotation
import C4Direction
import C4Replacement
import Mathlib.LinearAlgebra.Span.Basic

open scoped InnerProductSpace

namespace C4MaximalSparse

open C4Configuration

variable {E : Type*} [NormedAddCommGroup E] [InnerProductSpace ℝ E]

omit [InnerProductSpace ℝ E] in
theorem unit_norm (x : UnitVector E) : ‖x.1‖ = 1 := by
  simpa only [Metric.mem_sphere, dist_zero_right] using x.2

noncomputable def orthogonalNeighborSpan {n : ℕ} (v : Configuration E n) (i : Fin n) :
    Submodule ℝ E :=
  Submodule.span ℝ (Set.range (fun j : {j : Fin n // ⟪(v i).1, (v j).1⟫_ℝ = 0} =>
    (v j.1).1))

theorem neighbor_mem_span {n : ℕ} (v : Configuration E n) (i j : Fin n)
    (h : ⟪(v i).1, (v j).1⟫_ℝ = 0) :
    (v j).1 ∈ orthogonalNeighborSpan v i := by
  apply Submodule.subset_span
  exact ⟨⟨j, h⟩, rfl⟩

private theorem replacement_incident {n : ℕ} (v : Configuration E n) (i : Fin n)
    (w : UnitVector E) :
    C4Replacement.incidentAngle (C4Replacement.replacement v i w) i =
      ∑ j : {j : Fin n // j ≠ i}, Real.arccos |⟪w.1, (v j.1).1⟫_ℝ| := by
  classical
  rw [C4Replacement.incidentAngle_eq_sum_neighbors]
  apply Finset.sum_congr rfl
  intro j _
  simp only [C4Replacement.replacement, Function.update_self,
    Function.update_of_ne j.property]

/-- At a maximizer with the largest orthogonality count, a vertex that has a
nonorthogonal neighbor has an orthogonal-neighbor span of codimension at most one. -/
theorem orthogonalNeighbor_finrank [FiniteDimensional ℝ E] {n : ℕ}
    (v : Configuration E n) (i : Fin n)
    (hmax : ∀ w : Configuration E n, totalAngle w ≤ totalAngle v)
    (hselect : ∀ w : Configuration E n, totalAngle w = totalAngle v →
      orthogonalityCount w ≤ orthogonalityCount v)
    (hactive : ∃ j : Fin n, j ≠ i ∧ ⟪(v i).1, (v j).1⟫_ℝ ≠ 0) :
    Module.finrank ℝ E ≤ Module.finrank ℝ (orthogonalNeighborSpan v i) + 1 := by
  classical
  by_contra hn
  have hdim : Module.finrank ℝ (orthogonalNeighborSpan v i) + 1 <
      Module.finrank ℝ E := lt_of_not_ge hn
  obtain ⟨y, hy, hxy, hyW⟩ := C4Direction.exists_unit_orthogonal
    (orthogonalNeighborSpan v i) (v i).1 hdim
  let z : {j : Fin n // j ≠ i} → E := fun j => (v j.1).1
  have hz : ∀ j, ‖z j‖ = 1 := fun j => unit_norm (v j.1)
  have hactive' : ∃ j, ⟪(v i).1, z j⟫_ℝ ≠ 0 := by
    obtain ⟨j, hji, hj⟩ := hactive
    exact ⟨⟨j, hji⟩, hj⟩
  have hpreserve : ∀ j, ⟪(v i).1, z j⟫_ℝ = 0 → ⟪y, z j⟫_ℝ = 0 := by
    intro j hj
    exact hyW (z j) (neighbor_mem_span v i j.1 hj)
  have hstarmax : ∀ t : ℝ,
      (∑ j, Real.arccos |⟪C4Geometry.rotation (v i).1 y t, z j⟫_ℝ|) ≤
        ∑ j, Real.arccos |⟪(v i).1, z j⟫_ℝ| := by
    intro t
    let w : UnitVector E := ⟨C4Geometry.rotation (v i).1 y t, by
      simpa only [Metric.mem_sphere, dist_zero_right] using
        C4Geometry.rotation_unit (v i).1 y t (unit_norm (v i)) hy hxy⟩
    have hb := C4Replacement.incidentAngle_le_of_totalAngle_max v i w hmax
    rw [replacement_incident, C4Replacement.incidentAngle_eq_sum_neighbors] at hb
    exact hb
  obtain ⟨t, hunit, ht, hold, j, hj, hnew⟩ := C4LocalRotation.add_orthogonality
    (v i).1 y z (unit_norm (v i)) hy hxy hz hactive' hpreserve hstarmax
  let w : UnitVector E := ⟨C4Geometry.rotation (v i).1 y t, by
    simpa only [Metric.mem_sphere, dist_zero_right] using hunit⟩
  have hincident : C4Replacement.incidentAngle (C4Replacement.replacement v i w) i =
      C4Replacement.incidentAngle v i := by
    rw [replacement_incident, C4Replacement.incidentAngle_eq_sum_neighbors]
    exact ht
  have hscore : totalAngle (C4Replacement.replacement v i w) = totalAngle v :=
    C4Replacement.totalAngle_replacement_eq_of_incident_eq v i w hincident
  have hstrict : orthogonalityCount v <
      orthogonalityCount (C4Replacement.replacement v i w) := by
    apply C4Replacement.orthogonalityCount_replacement_strict v i w
    · intro k hki hk
      exact hold ⟨k, hki⟩ hk
    · exact ⟨j.1, j.2, hj, hnew⟩
  exact (not_lt_of_ge (hselect (C4Replacement.replacement v i w) hscore)) hstrict

#print axioms neighbor_mem_span
#print axioms orthogonalNeighbor_finrank
end C4MaximalSparse
\end{Verbatim}

\subsubsection{14. C4NeighborCount.lean}\label{c4neighborcount.lean}

\href{https://github.com/mpelteshki/climbing-to-the-frontier/blob/29fbb70da18564048f62492e43870f2f96a65463/cagent/p1/C4NeighborCount.lean}{Pinned source file} · SHA-256 \texttt{9d4f9fa2\allowbreak{}f09f8051\allowbreak{}92d2ac8b\allowbreak{}9d9a3e98\allowbreak{}01615b11\allowbreak{}b0f513fa\allowbreak{}83e0c1ab\allowbreak{}b5d74833}

\begin{Verbatim}[breaklines,breakanywhere,fontsize=\footnotesize]
import C4MaximalSparse
import Mathlib.LinearAlgebra.Dimension.Constructions
import Mathlib.Data.Finset.Card

open scoped InnerProductSpace

namespace C4NeighborCount

open C4Configuration C4MaximalSparse

variable {E : Type*} [NormedAddCommGroup E] [InnerProductSpace ℝ E]

noncomputable def nonorthogonalNeighbors {n : ℕ}
    (v : Configuration E n) (i : Fin n) : Finset (Fin n) :=
  Finset.univ.filter (fun j => j ≠ i ∧ ⟪(v i).1, (v j).1⟫_ℝ ≠ 0)

theorem neighbor_span_finrank_le_zero_count {n : ℕ}
    (v : Configuration E n) (i : Fin n) :
    Module.finrank ℝ (orthogonalNeighborSpan v i) ≤
      (Finset.univ.filter (fun j : Fin n =>
        ⟪(v i).1, (v j).1⟫_ℝ = 0)).card := by
  classical
  have hcard : Fintype.card {j : Fin n // ⟪(v i).1, (v j).1⟫_ℝ = 0} =
      (Finset.univ.filter (fun j : Fin n =>
        ⟪(v i).1, (v j).1⟫_ℝ = 0)).card := by
    rw [Fintype.card_subtype]
  have h := finrank_range_le_card (R := ℝ)
    (fun j : {j : Fin n // ⟪(v i).1, (v j).1⟫_ℝ = 0} => (v j.1).1)
  rw [hcard] at h
  convert h using 1
  rfl

theorem nonorthogonal_neighbors_le {n d : ℕ} (v : Configuration E n)
    (hdn : d ≤ n)
    (hrank : ∀ i : Fin n, (∃ j : Fin n,
        j ≠ i ∧ ⟪(v i).1, (v j).1⟫_ℝ ≠ 0) →
      d ≤ Module.finrank ℝ (orthogonalNeighborSpan v i) + 1)
    (i : Fin n) : (nonorthogonalNeighbors v i).card ≤ n - d := by
  classical
  let P : Fin n → Prop := fun j => ⟪(v i).1, (v j).1⟫_ℝ = 0
  let Z := Finset.univ.filter P
  let N := Finset.univ.filter (fun j => ¬ P j)
  have hself : ¬ P i := by
    dsimp [P]
    rw [real_inner_self_eq_norm_sq, unit_norm (v i)]
    norm_num
  have hiN : i ∈ N := by simp [N, hself]
  have hdegree : (nonorthogonalNeighbors v i).card + 1 = N.card := by
    have hset : nonorthogonalNeighbors v i = N.erase i := by
      ext j
      simp [nonorthogonalNeighbors, N, P, Finset.mem_erase]
    rw [hset]
    exact Finset.card_erase_add_one hiN
  have hpartition : Z.card + N.card = n := by
    simpa [Z, N, P] using
      (Finset.card_filter_add_card_filter_not (s := Finset.univ) P)
  by_cases hactive : ∃ j : Fin n,
      j ≠ i ∧ ⟪(v i).1, (v j).1⟫_ℝ ≠ 0
  · have hspan := neighbor_span_finrank_le_zero_count v i
    have hlower := hrank i hactive
    change Module.finrank ℝ (orthogonalNeighborSpan v i) ≤ Z.card at hspan
    omega
  · have hempty : nonorthogonalNeighbors v i = ∅ := by
      apply Finset.eq_empty_iff_forall_notMem.mpr
      intro j hj
      exact hactive ⟨j, (Finset.mem_filter.mp hj).2⟩
    simp [hempty]

#print axioms C4NeighborCount.neighbor_span_finrank_le_zero_count
#print axioms C4NeighborCount.nonorthogonal_neighbors_le

end C4NeighborCount
\end{Verbatim}

\subsubsection{15. C4SparseConfiguration.lean}\label{c4sparseconfiguration.lean}

\href{https://github.com/mpelteshki/climbing-to-the-frontier/blob/29fbb70da18564048f62492e43870f2f96a65463/cagent/p1/C4SparseConfiguration.lean}{Pinned source file} · SHA-256 \texttt{7bd4a7aa\allowbreak{}9e3b8b48\allowbreak{}313688d6\allowbreak{}e6126ab2\allowbreak{}19cf60c5\allowbreak{}4ad27508\allowbreak{}9fe8e9fd\allowbreak{}050f9932}

\begin{Verbatim}[breaklines,breakanywhere,fontsize=\footnotesize]
import C4NeighborCount

namespace C4SparseConfiguration

open C4Configuration

variable {E : Type*} [NormedAddCommGroup E] [InnerProductSpace ℝ E]
  [FiniteDimensional ℝ E] [ProperSpace E]

/-- An angle-maximizing configuration can be chosen so that each line has at
most `n - dim E` nonorthogonal neighbors. -/
theorem exists_sparse_maximizer {n : ℕ} (hdn : Module.finrank ℝ E ≤ n)
    (u : UnitVector E) :
    ∃ v : Configuration E n,
      (∀ w : Configuration E n, totalAngle w ≤ totalAngle v) ∧
      (∀ w : Configuration E n, totalAngle w = totalAngle v →
        orthogonalityCount w ≤ orthogonalityCount v) ∧
      ∀ i : Fin n, (C4NeighborCount.nonorthogonalNeighbors v i).card ≤
        n - Module.finrank ℝ E := by
  obtain ⟨v, hmax, hselect⟩ := exists_lex_max (n := n) u
  refine ⟨v, hmax, hselect, ?_⟩
  exact C4NeighborCount.nonorthogonal_neighbors_le v hdn
    (fun i hi => C4MaximalSparse.orthogonalNeighbor_finrank v i hmax hselect hi)

/-- The C4 reduction to a nonorthogonality graph of maximum degree two. -/
theorem exists_degree_two_maximizer (u : UnitVector E) :
    ∃ v : Configuration E (Module.finrank ℝ E + 2),
      (∀ w : Configuration E (Module.finrank ℝ E + 2), totalAngle w ≤ totalAngle v) ∧
      ∀ i, (C4NeighborCount.nonorthogonalNeighbors v i).card ≤ 2 := by
  obtain ⟨v, hmax, _, hdegree⟩ :=
    exists_sparse_maximizer (n := Module.finrank ℝ E + 2) (by omega) u
  refine ⟨v, hmax, ?_⟩
  simpa using hdegree

#print axioms exists_sparse_maximizer
#print axioms exists_degree_two_maximizer
end C4SparseConfiguration
\end{Verbatim}

\subsubsection{16. C4MatrixCorank.lean}\label{c4matrixcorank.lean}

\href{https://github.com/mpelteshki/climbing-to-the-frontier/blob/29fbb70da18564048f62492e43870f2f96a65463/cagent/p1/C4MatrixCorank.lean}{Pinned source file} · SHA-256 \texttt{180d0072\allowbreak{}3bd73798\allowbreak{}394275b5\allowbreak{}9b01dfa8\allowbreak{}5a2c797b\allowbreak{}335e8267\allowbreak{}dd121993\allowbreak{}b5224cee}

\begin{Verbatim}[breaklines,breakanywhere,fontsize=\footnotesize]
import Mathlib.LinearAlgebra.Matrix.Rank
import Mathlib.LinearAlgebra.FiniteDimensional.Lemmas
import Mathlib.LinearAlgebra.Dimension.StrongRankCondition
import Mathlib.Analysis.SpecialFunctions.Trigonometric.Inverse

/-! Corank bounds for matrices with a nonzero interior superdiagonal. -/
namespace C4MatrixCorank

private theorem kernel_zero_of_first_two {n : ℕ} (hn : 2 ≤ n)
    (M : Matrix (Fin n) (Fin n) ℝ)
    (hband : ∀ i j : Fin n, 1 ≤ i.val → i.val + 1 < n →
      (j.val + 1 < i.val ∨ i.val + 1 < j.val) → M i j = 0)
    (hsuper : ∀ i j : Fin n, 1 ≤ i.val → i.val + 1 < n →
      j.val = i.val + 1 → M i j ≠ 0)
    (x : Fin n → ℝ) (hMx : M.mulVecLin x = 0)
    (h0 : x ⟨0, by omega⟩ = 0) (h1 : x ⟨1, by omega⟩ = 0) : x = 0 := by
  funext j
  have hcoord : ∀ k : ℕ, (hk : k < n) → x ⟨k, hk⟩ = 0 := by
    intro k
    induction k using Nat.strong_induction_on with
    | h k ih =>
      intro hk
      by_cases hk0 : k = 0
      · subst k
        simpa using h0
      by_cases hk1 : k = 1
      · subst k
        simpa using h1
      have hk2 : 2 ≤ k := by omega
      let i : Fin n := ⟨k - 1, by omega⟩
      let target : Fin n := ⟨k, hk⟩
      have hi1 : 1 ≤ i.val := by dsimp [i]; omega
      have hi2 : i.val + 1 < n := by dsimp [i]; omega
      have hrow : ∑ t : Fin n, M i t * x t = 0 := by
        have hh := congrArg (fun f : Fin n → ℝ => f i) hMx
        simpa only [Matrix.mulVecLin_apply, Matrix.mulVec_apply_eq_sum,
          Pi.zero_apply] using hh
      have hsum : (∑ t : Fin n, M i t * x t) = M i target * x target := by
        classical
        apply Finset.sum_eq_single target
        · intro t _ hne
          by_cases htk : t.val < k
          · have ht0 : x t = 0 := by simpa using ih t.val htk t.isLt
            simp [ht0]
          · have htail : i.val + 1 < t.val := by
              dsimp [i]
              have : t.val ≠ k := by
                intro heq
                apply hne
                exact Fin.ext heq
              omega
            simp [hband i t hi1 hi2 (Or.inr htail)]
        · simp
      have htarget : M i target * x target = 0 := by rw [← hsum]; exact hrow
      exact (mul_eq_zero.mp htarget).resolve_left
        (hsuper i target hi1 hi2 (by dsimp [i, target]; omega))
  simpa using hcoord j.val j.isLt

/-- At most two independent kernel vectors can survive the interior recurrence. -/
theorem finrank_ker_le_two {n : ℕ} (hn : 2 ≤ n)
    (M : Matrix (Fin n) (Fin n) ℝ)
    (hband : ∀ i j : Fin n, 1 ≤ i.val → i.val + 1 < n →
      (j.val + 1 < i.val ∨ i.val + 1 < j.val) → M i j = 0)
    (hsuper : ∀ i j : Fin n, 1 ≤ i.val → i.val + 1 < n →
      j.val = i.val + 1 → M i j ≠ 0) :
    Module.finrank ℝ (LinearMap.ker M.mulVecLin) ≤ 2 := by
  let K := LinearMap.ker M.mulVecLin
  let obs : K →ₗ[ℝ] ℝ × ℝ := {
    toFun := fun x => (x.1 ⟨0, by omega⟩, x.1 ⟨1, by omega⟩)
    map_add' := by intro x y; rfl
    map_smul' := by intro r x; rfl }
  have hinj : Function.Injective obs := by
    intro x y hxy
    have h0 : (x - y).1 ⟨0, by omega⟩ = 0 := by
      have h := congrArg Prod.fst hxy
      change x.1 ⟨0, by omega⟩ = y.1 ⟨0, by omega⟩ at h
      simpa using sub_eq_zero.mpr h
    have h1 : (x - y).1 ⟨1, by omega⟩ = 0 := by
      have h := congrArg Prod.snd hxy
      change x.1 ⟨1, by omega⟩ = y.1 ⟨1, by omega⟩ at h
      simpa using sub_eq_zero.mpr h
    have hz : x - y = 0 := by
      apply Subtype.ext
      exact kernel_zero_of_first_two hn M hband hsuper (x - y).1 (x - y).2 h0 h1
    exact sub_eq_zero.mp hz
  simpa using obs.finrank_le_finrank_of_injective hinj

/-- The range of a cycle-tridiagonal matrix has codimension at most two. -/
theorem finrank_range_add_two_ge {n : ℕ} (hn : 2 ≤ n)
    (M : Matrix (Fin n) (Fin n) ℝ)
    (hband : ∀ i j : Fin n, 1 ≤ i.val → i.val + 1 < n →
      (j.val + 1 < i.val ∨ i.val + 1 < j.val) → M i j = 0)
    (hsuper : ∀ i j : Fin n, 1 ≤ i.val → i.val + 1 < n →
      j.val = i.val + 1 → M i j ≠ 0) :
    n ≤ Module.finrank ℝ (LinearMap.range M.mulVecLin) + 2 := by
  have hker := finrank_ker_le_two hn M hband hsuper
  have hnull := (M.mulVecLin).finrank_range_add_finrank_ker
  have hdim : Module.finrank ℝ (Fin n → ℝ) = n := by simp
  omega

/-- A cycle-tridiagonal matrix with nonzero interior superdiagonal has corank at most two. -/
theorem rank_add_two_ge {n : ℕ} (hn : 2 ≤ n)
    (M : Matrix (Fin n) (Fin n) ℝ)
    (hband : ∀ i j : Fin n, 1 ≤ i.val → i.val + 1 < n →
      (j.val + 1 < i.val ∨ i.val + 1 < j.val) → M i j = 0)
    (hsuper : ∀ i j : Fin n, 1 ≤ i.val → i.val + 1 < n →
      j.val = i.val + 1 → M i j ≠ 0) :
    n ≤ M.rank + 2 := by
  exact finrank_range_add_two_ge hn M hband hsuper

private theorem path_kernel_zero_of_first_one {n : ℕ} (hn : 2 ≤ n)
    (M : Matrix (Fin n) (Fin n) ℝ)
    (hband : ∀ i j : Fin n, 1 ≤ i.val → i.val + 1 < n →
      (j.val + 1 < i.val ∨ i.val + 1 < j.val) → M i j = 0)
    (hsuper : ∀ i j : Fin n, 1 ≤ i.val → i.val + 1 < n →
      j.val = i.val + 1 → M i j ≠ 0)
    (hfirst : ∀ j : Fin n, 1 < j.val → M ⟨0, by omega⟩ j = 0)
    (h01 : M ⟨0, by omega⟩ ⟨1, by omega⟩ ≠ 0)
    (x : Fin n → ℝ) (hMx : M.mulVecLin x = 0)
    (h0 : x ⟨0, by omega⟩ = 0) : x = 0 := by
  let i0 : Fin n := ⟨0, by omega⟩
  let j1 : Fin n := ⟨1, by omega⟩
  have hrow : ∑ j : Fin n, M i0 j * x j = 0 := by
    have hh := congrArg (fun f : Fin n → ℝ => f i0) hMx
    simpa only [Matrix.mulVecLin_apply, Matrix.mulVec_apply_eq_sum,
      Pi.zero_apply] using hh
  have hsum : (∑ j : Fin n, M i0 j * x j) = M i0 j1 * x j1 := by
    classical
    apply Finset.sum_eq_single j1
    · intro j _ hne
      by_cases hj0 : j.val = 0
      · have : j = i0 := Fin.ext hj0
        subst j
        simp [i0, h0]
      · have hjgt : 1 < j.val := by
          have hj1 : j.val ≠ 1 := by
            intro heq
            apply hne
            exact Fin.ext heq
          omega
        have hz : M i0 j = 0 := by simpa only [i0] using hfirst j hjgt
        simp [hz]
    · simp
  have h1 : x j1 = 0 := by
    have hprod : M i0 j1 * x j1 = 0 := by rw [← hsum]; exact hrow
    exact (mul_eq_zero.mp hprod).resolve_left (by simpa [i0, j1] using h01)
  exact kernel_zero_of_first_two hn M hband hsuper x hMx h0 (by simpa [j1] using h1)

/-- A path matrix with a nonzero first superdiagonal entry has corank at most one. -/
theorem path_finrank_ker_le_one {n : ℕ} (hn : 2 ≤ n)
    (M : Matrix (Fin n) (Fin n) ℝ)
    (hband : ∀ i j : Fin n, 1 ≤ i.val → i.val + 1 < n →
      (j.val + 1 < i.val ∨ i.val + 1 < j.val) → M i j = 0)
    (hsuper : ∀ i j : Fin n, 1 ≤ i.val → i.val + 1 < n →
      j.val = i.val + 1 → M i j ≠ 0)
    (hfirst : ∀ j : Fin n, 1 < j.val → M ⟨0, by omega⟩ j = 0)
    (h01 : M ⟨0, by omega⟩ ⟨1, by omega⟩ ≠ 0) :
    Module.finrank ℝ (LinearMap.ker M.mulVecLin) ≤ 1 := by
  let K := LinearMap.ker M.mulVecLin
  let obs : K →ₗ[ℝ] ℝ := {
    toFun := fun x => x.1 ⟨0, by omega⟩
    map_add' := by intro x y; rfl
    map_smul' := by intro r x; rfl }
  have hinj : Function.Injective obs := by
    intro x y hxy
    have h0 : (x - y).1 ⟨0, by omega⟩ = 0 := by
      change x.1 ⟨0, by omega⟩ = y.1 ⟨0, by omega⟩ at hxy
      simpa using sub_eq_zero.mpr hxy
    have hz : x - y = 0 := by
      apply Subtype.ext
      exact path_kernel_zero_of_first_one hn M hband hsuper hfirst h01
        (x - y).1 (x - y).2 h0
    exact sub_eq_zero.mp hz
  simpa using obs.finrank_le_finrank_of_injective hinj

/-- A path matrix with the stated band and nonzero superdiagonal has rank at least `n-1`. -/
theorem path_rank_add_one_ge {n : ℕ} (hn : 2 ≤ n)
    (M : Matrix (Fin n) (Fin n) ℝ)
    (hband : ∀ i j : Fin n, 1 ≤ i.val → i.val + 1 < n →
      (j.val + 1 < i.val ∨ i.val + 1 < j.val) → M i j = 0)
    (hsuper : ∀ i j : Fin n, 1 ≤ i.val → i.val + 1 < n →
      j.val = i.val + 1 → M i j ≠ 0)
    (hfirst : ∀ j : Fin n, 1 < j.val → M ⟨0, by omega⟩ j = 0)
    (h01 : M ⟨0, by omega⟩ ⟨1, by omega⟩ ≠ 0) :
    n ≤ M.rank + 1 := by
  have hker := path_finrank_ker_le_one hn M hband hsuper hfirst h01
  have hnull := (M.mulVecLin).finrank_range_add_finrank_ker
  have hdim : Module.finrank ℝ (Fin n → ℝ) = n := by simp
  change M.rank + Module.finrank ℝ (LinearMap.ker M.mulVecLin) =
    Module.finrank ℝ (Fin n → ℝ) at hnull
  omega

#print axioms finrank_ker_le_two
#print axioms finrank_range_add_two_ge
#print axioms rank_add_two_ge
#print axioms path_finrank_ker_le_one
#print axioms path_rank_add_one_ge
end C4MatrixCorank
\end{Verbatim}

\subsubsection{17. C4FiveCycle.lean}\label{c4fivecycle.lean}

\href{https://github.com/mpelteshki/climbing-to-the-frontier/blob/29fbb70da18564048f62492e43870f2f96a65463/cagent/p1/C4FiveCycle.lean}{Pinned source file} · SHA-256 \texttt{fe0bb40d\allowbreak{}9fc9b71f\allowbreak{}42d69ef8\allowbreak{}14c7199d\allowbreak{}b5a214bf\allowbreak{}bcee507a\allowbreak{}ddca3340\allowbreak{}f8d4aed7}

\begin{Verbatim}[breaklines,breakanywhere,fontsize=\footnotesize]
import C4Pentagon

/-! The explicit coordinate five-cycle scalar inequality. -/
namespace C4FiveCycle

private theorem asin_ratio_data {x y L : ℝ} (hx0 : 0 ≤ x) (hy0 : 0 ≤ y)
    (hLpos : 0 < L) (hL2 : L ^ 2 = x ^ 2 + y ^ 2) :
    Real.sin (Real.arcsin (x / L)) = x / L ∧
      Real.cos (Real.arcsin (x / L)) = y / L ∧
      0 ≤ Real.arcsin (x / L) ∧ Real.arcsin (x / L) ≤ Real.pi / 2 := by
  have hxL : x ≤ L := by nlinarith only [hL2, hx0, hy0, hLpos, sq_nonneg y]
  have hxratio0 : 0 ≤ x / L := div_nonneg hx0 hLpos.le
  have hxratio1 : x / L ≤ 1 := (div_le_one hLpos).mpr hxL
  have hcos : Real.cos (Real.arcsin (x / L)) = y / L := by
    rw [Real.cos_arcsin]
    have hrad : 0 ≤ 1 - (x / L) ^ 2 := by nlinarith [sq_nonneg (x / L - 1)]
    have hsq : (Real.sqrt (1 - (x / L) ^ 2)) ^ 2 = (y / L) ^ 2 := by
      rw [Real.sq_sqrt hrad]
      calc
        1 - (x / L) ^ 2 = (L ^ 2 - x ^ 2) / L ^ 2 := by field_simp
        _ = (y / L) ^ 2 := by rw [div_pow, hL2]; ring
    nlinarith only [hsq, Real.sqrt_nonneg (1 - (x / L) ^ 2),
      div_nonneg hy0 hLpos.le]
  exact ⟨Real.sin_arcsin (by linarith) hxratio1, hcos,
    Real.arcsin_nonneg.mpr hxratio0, Real.arcsin_le_pi_div_two _⟩

set_option maxHeartbeats 1000000 in
/-- Five-cycle deficiency bound for the nondegenerate coordinate model. -/
theorem explicit_five_cycle {a c : ℝ} (ha0 : 0 < a) (ha1 : a < 1)
    (hc0 : 0 < c) (hc1 : c < 1) :
    Real.pi ≤
      Real.arcsin (c * Real.sqrt (1 - a ^ 2) /
        Real.sqrt (a ^ 2 + c ^ 2 - a ^ 2 * c ^ 2)) +
      Real.arcsin (a * Real.sqrt (1 - c ^ 2) /
        Real.sqrt (a ^ 2 + c ^ 2 - a ^ 2 * c ^ 2)) +
      Real.arcsin c +
      Real.arcsin (Real.sqrt (1 - a ^ 2) * Real.sqrt (1 - c ^ 2)) +
      Real.arcsin a := by
  let s := Real.sqrt (1 - a ^ 2)
  let t := Real.sqrt (1 - c ^ 2)
  let L := Real.sqrt (a ^ 2 + c ^ 2 - a ^ 2 * c ^ 2)
  let u := c * s / L
  let v := a * t / L
  let E := Real.arcsin (a * c / (1 + s * t))
  have hs0 : 0 ≤ s := Real.sqrt_nonneg _
  have ht0 : 0 ≤ t := Real.sqrt_nonneg _
  have hs2 : s ^ 2 = 1 - a ^ 2 := Real.sq_sqrt (by nlinarith)
  have ht2 : t ^ 2 = 1 - c ^ 2 := Real.sq_sqrt (by nlinarith)
  have hs1 : s ≤ 1 := by nlinarith only [hs0, hs2, sq_nonneg a]
  have ht1 : t ≤ 1 := by nlinarith only [ht0, ht2, sq_nonneg c]
  have hL2 : L ^ 2 = a ^ 2 + c ^ 2 - a ^ 2 * c ^ 2 := by
    apply Real.sq_sqrt
    nlinarith [mul_nonneg (sq_nonneg c) (by nlinarith only [ha0.le, ha1.le])]
  have hLpos : 0 < L := by
    have hL0 := Real.sqrt_nonneg (a ^ 2 + c ^ 2 - a ^ 2 * c ^ 2)
    have haux := mul_nonneg (sq_nonneg c) (by nlinarith only [ha0.le, ha1.le] : 0 ≤ 1 - a ^ 2)
    dsimp [L] at *
    nlinarith only [hL0, hL2, haux, ha0]
  have hxsq : L ^ 2 = (c * s) ^ 2 + a ^ 2 := by
    rw [mul_pow, hs2, hL2]
    ring
  have hysq : L ^ 2 = (a * t) ^ 2 + c ^ 2 := by
    rw [mul_pow, ht2, hL2]
    ring
  have hu := asin_ratio_data (mul_nonneg hc0.le hs0) ha0.le hLpos hxsq
  have hv := asin_ratio_data (mul_nonneg ha0.le ht0) hc0.le hLpos hysq
  have hst0 : 0 ≤ s * t := mul_nonneg hs0 ht0
  have hst1 : s * t ≤ 1 := by nlinarith [mul_nonneg (sub_nonneg.mpr hs1) ht0]
  have hdpos : 0 < 1 + s * t := by linarith
  have hprod : (s * t) ^ 2 = (1 - a ^ 2) * (1 - c ^ 2) := by
    rw [mul_pow, hs2, ht2]
  have hLrelation : L ^ 2 = 1 - (s * t) ^ 2 := by
    rw [hL2, hprod]
    ring
  have hratioE0 : 0 ≤ a * c / (1 + s * t) := by positivity
  have hratioE1 : a * c / (1 + s * t) ≤ 1 := by
    apply (div_le_one hdpos).mpr
    nlinarith [mul_nonneg (sub_nonneg.mpr ha1.le) hc0.le]
  have hsinE : Real.sin E = a * c / (1 + s * t) :=
    Real.sin_arcsin (by linarith) hratioE1
  have hE0 : 0 ≤ E := Real.arcsin_nonneg.mpr hratioE0
  have hE1 : E ≤ Real.pi / 2 := Real.arcsin_le_pi_div_two _
  have hsumcos : Real.cos (Real.arcsin u + Real.arcsin v) =
      a * c / (1 + s * t) := by
    rw [Real.cos_add, hu.2.1, hv.2.1, hu.1, hv.1]
    calc
      a / L * (c / L) - c * s / L * (a * t / L) =
          (a * c * (1 - s * t)) / L ^ 2 := by ring
      _ = a * c / (1 + s * t) := by
        have hLne : L ^ 2 ≠ 0 := pow_ne_zero 2 (ne_of_gt hLpos)
        field_simp
        nlinarith only [hLrelation]
  have hsum0 : 0 ≤ Real.arcsin u + Real.arcsin v := add_nonneg hu.2.2.1 hv.2.2.1
  have hsumpi : Real.arcsin u + Real.arcsin v ≤ Real.pi := by
    linarith [hu.2.2.2, hv.2.2.2]
  have hsum : Real.arcsin u + Real.arcsin v = Real.pi / 2 - E := by
    have hcosTarget : Real.cos (Real.pi / 2 - E) = a * c / (1 + s * t) := by
      rw [Real.cos_pi_div_two_sub, hsinE]
    have htarget0 : 0 ≤ Real.pi / 2 - E := by linarith
    have htargetpi : Real.pi / 2 - E ≤ Real.pi := by linarith [Real.pi_pos]
    have h1 := Real.arccos_cos hsum0 hsumpi
    have h2 := Real.arccos_cos htarget0 htargetpi
    rw [hsumcos] at h1
    rw [hcosTarget] at h2
    linarith
  have hacosS : Real.arccos s = Real.arcsin a := by
    have haasin0 : 0 ≤ Real.arcsin a := Real.arcsin_nonneg.mpr ha0.le
    have haasin1 : Real.arcsin a ≤ Real.pi := by
      linarith [Real.arcsin_le_pi_div_two a, Real.pi_pos]
    rw [← Real.arccos_cos haasin0 haasin1, Real.cos_arcsin]
  have hacosT : Real.arccos t = Real.arcsin c := by
    have hcasin0 : 0 ≤ Real.arcsin c := Real.arcsin_nonneg.mpr hc0.le
    have hcasin1 : Real.arcsin c ≤ Real.pi := by
      linarith [Real.arcsin_le_pi_div_two c, Real.pi_pos]
    rw [← Real.arccos_cos hcasin0 hcasin1, Real.cos_arcsin]
  have hP := C4Pentagon.pentagon_scalar hs0 hs1 ht0 ht1
  rw [show 1 - s ^ 2 = a ^ 2 by nlinarith [hs2],
      show 1 - t ^ 2 = c ^ 2 by nlinarith [ht2],
      Real.sqrt_sq ha0.le, Real.sqrt_sq hc0.le, hacosS, hacosT] at hP
  have hst : Real.arcsin (s * t) = Real.pi / 2 - Real.arccos (s * t) :=
    Real.arcsin_eq_pi_div_two_sub_arccos _
  change Real.arccos (s * t) + E ≤ Real.arcsin a + Real.arcsin c at hP
  change Real.pi ≤ Real.arcsin u + Real.arcsin v + Real.arcsin c +
    Real.arcsin (s * t) + Real.arcsin a
  linarith only [hsum, hP, hst]

#print axioms explicit_five_cycle
end C4FiveCycle
\end{Verbatim}

\subsubsection{18. C4SixCycle.lean}\label{c4sixcycle.lean}

\href{https://github.com/mpelteshki/climbing-to-the-frontier/blob/29fbb70da18564048f62492e43870f2f96a65463/cagent/p1/C4SixCycle.lean}{Pinned source file} · SHA-256 \texttt{2dc09cb3\allowbreak{}ac16e054\allowbreak{}a07e888b\allowbreak{}0e9aa527\allowbreak{}f541f453\allowbreak{}7c84b906\allowbreak{}b34ec064\allowbreak{}6df1be6f}

\begin{Verbatim}[breaklines,breakanywhere,fontsize=\footnotesize]
import C4Pentagon

/-! The local six-cycle projection inequality in raw coordinates. -/
namespace C4SixCycle

private theorem arccos_ratio_data {x y L : ℝ} (hx0 : 0 ≤ x) (hy0 : 0 ≤ y)
    (hLpos : 0 < L) (hL2 : L ^ 2 = x ^ 2 + y ^ 2) :
    Real.cos (Real.arccos (x / L)) = x / L ∧
      Real.sin (Real.arccos (x / L)) = y / L ∧
      0 ≤ Real.arccos (x / L) ∧ Real.arccos (x / L) ≤ Real.pi / 2 := by
  have hxL : x ≤ L := by nlinarith only [hL2, hx0, hy0, hLpos, sq_nonneg y]
  have hxratio0 : 0 ≤ x / L := div_nonneg hx0 hLpos.le
  have hxratio1 : x / L ≤ 1 := (div_le_one hLpos).mpr hxL
  have hsin : Real.sin (Real.arccos (x / L)) = y / L := by
    rw [Real.sin_arccos]
    have hrad : 0 ≤ 1 - (x / L) ^ 2 := by nlinarith [sq_nonneg (x / L - 1)]
    have hsq : (Real.sqrt (1 - (x / L) ^ 2)) ^ 2 = (y / L) ^ 2 := by
      rw [Real.sq_sqrt hrad]
      calc
        1 - (x / L) ^ 2 = (L ^ 2 - x ^ 2) / L ^ 2 := by field_simp
        _ = (y / L) ^ 2 := by rw [div_pow, hL2]; ring
    nlinarith only [hsq, Real.sqrt_nonneg (1 - (x / L) ^ 2),
      div_nonneg hy0 hLpos.le]
  exact ⟨Real.cos_arccos (by linarith) hxratio1, hsin,
    Real.arccos_nonneg _, Real.arccos_le_pi_div_two.mpr hxratio0⟩

/-- Raw-coordinate local projection estimate for an irreducible six-cycle. -/
theorem raw_projection_inequality {a b c d u v : ℝ}
    (ha0 : 0 < a) (hb0 : 0 < b) (hc0 : 0 < c) (hd0 : 0 < d)
    (hu0 : 0 < u) (hv0 : 0 < v)
    (hleft : a ^ 2 + b ^ 2 + u ^ 2 = 1)
    (hright : c ^ 2 + d ^ 2 + v ^ 2 = 1)
    (hcross : u * v = b * c) :
    Real.arcsin (a / Real.sqrt (1 - b ^ 2)) +
      Real.arcsin (d / Real.sqrt (1 - c ^ 2)) +
      Real.arcsin (b * c /
        (Real.sqrt (1 - b ^ 2) * Real.sqrt (1 - c ^ 2))) ≤
      Real.arcsin a + Real.arcsin b + Real.arcsin c + Real.arcsin d := by
  let A := Real.sqrt (1 - b ^ 2)
  let D := Real.sqrt (1 - c ^ 2)
  let R := Real.sqrt (b ^ 2 + u ^ 2)
  let B := Real.arcsin b
  let C := Real.arcsin c
  let U := Real.arccos (a / A)
  let V := Real.arccos (d / D)
  let T := Real.arccos (b / R)
  have hA2 : A ^ 2 = 1 - b ^ 2 := Real.sq_sqrt (by nlinarith only [hleft, sq_nonneg a, sq_nonneg u])
  have hD2 : D ^ 2 = 1 - c ^ 2 := Real.sq_sqrt (by nlinarith only [hright, sq_nonneg d, sq_nonneg v])
  have hR2 : R ^ 2 = b ^ 2 + u ^ 2 := Real.sq_sqrt (by positivity)
  have hA2' : A ^ 2 = a ^ 2 + u ^ 2 := by nlinarith only [hA2, hleft]
  have hD2' : D ^ 2 = d ^ 2 + v ^ 2 := by nlinarith only [hD2, hright]
  have hApos : 0 < A := by
    have hA0 := Real.sqrt_nonneg (1 - b ^ 2)
    nlinarith only [hA0, hA2', ha0, sq_nonneg u]
  have hDpos : 0 < D := by
    have hD0 := Real.sqrt_nonneg (1 - c ^ 2)
    nlinarith only [hD0, hD2', hd0, sq_nonneg v]
  have hRpos : 0 < R := by
    have hR0 := Real.sqrt_nonneg (b ^ 2 + u ^ 2)
    nlinarith only [hR0, hR2, hb0, sq_nonneg u]
  have hb1 : b ≤ 1 := by nlinarith only [hleft, hb0, ha0, hu0]
  have hc1 : c ≤ 1 := by nlinarith only [hright, hc0, hd0, hv0]
  have hBsin : Real.sin B = b := Real.sin_arcsin (by linarith) hb1
  have hCsin : Real.sin C = c := Real.sin_arcsin (by linarith) hc1
  have hBcos : Real.cos B = A := Real.cos_arcsin b
  have hCcos : Real.cos C = D := Real.cos_arcsin c
  have hB0 : 0 ≤ B := Real.arcsin_nonneg.mpr hb0.le
  have hC0 : 0 ≤ C := Real.arcsin_nonneg.mpr hc0.le
  have hB1 : B ≤ Real.pi / 2 := Real.arcsin_le_pi_div_two _
  have hC1 : C ≤ Real.pi / 2 := Real.arcsin_le_pi_div_two _
  have hU := arccos_ratio_data ha0.le hu0.le hApos hA2'
  have hV := arccos_ratio_data hd0.le hv0.le hDpos hD2'
  have hT := arccos_ratio_data hb0.le hu0.le hRpos hR2
  have hrel1 : Real.sin B * Real.sin T = Real.cos B * Real.sin U * Real.cos T := by
    rw [hBsin, hT.2.1, hBcos, hU.2.1, hT.1]
    field_simp
  have hrel2 : Real.sin C * Real.cos T = Real.cos C * Real.sin V * Real.sin T := by
    rw [hCsin, hT.1, hCcos, hV.2.1, hT.2.1]
    field_simp
    nlinarith only [hcross]
  have hproj := C4Pentagon.local_projection_inequality
    hT.2.2.1 hT.2.2.2 hB0 hB1 hC0 hC1
    hU.2.2.1 hU.2.2.2 hV.2.2.1 hV.2.2.2 hrel1 hrel2
  have hsinUV : Real.sin U * Real.sin V = b * c / (A * D) := by
    rw [hU.2.1, hV.2.1, ← hcross]
    ring
  have hcosBU : Real.cos B * Real.cos U = a := by
    rw [hBcos, hU.1]
    field_simp
  have hcosCV : Real.cos C * Real.cos V = d := by
    rw [hCcos, hV.1]
    field_simp
  rw [hsinUV, hcosBU, hcosCV] at hproj
  simp only [Real.arccos_eq_pi_div_two_sub_arcsin] at hproj
  change Real.arcsin (a / A) + Real.arcsin (d / D) +
    Real.arcsin (b * c / (A * D)) ≤
    Real.arcsin a + Real.arcsin b + Real.arcsin c + Real.arcsin d
  dsimp [B, C] at hproj
  linarith

#print axioms raw_projection_inequality
end C4SixCycle
\end{Verbatim}
\end{document}
